% numodel.dtx
%
% Docstrip source for numodel.sty, numodel-EN.def and numodel-NL.def.
% Run
%   tex numodel.ins
% to extract the derived files.  User-facing documentation lives
% in numodel-manual.tex (a stand-alone LaTeX file).
%
% Copyright (C) 2026 Paul Zuurbier <mail@paulzuurbier.nl>
%
% This work may be distributed and/or modified under the conditions
% of the LaTeX Project Public License, either version 1.3c of this
% license or (at your option) any later version.  The latest version
% of this license is in https://www.latex-project.org/lppl.txt
%
% This work has the LPPL maintenance status 'maintained'.
% The Current Maintainer of this work is Paul Zuurbier.
%
% This work consists of the files numodel.dtx and numodel.ins, the
% derived files numodel.sty, numodel-EN.def and numodel-NL.def, and the
% companion Lua module numodel.lua.
% \section{Implementation}
%
%    \begin{macrocode}
%<*package>
\NeedsTeXFormat{LaTeX2e}
\ProvidesPackage{numodel}[2026/05/26 v0.6.0
  Numerical physics models with Euler integration and Forrester diagrams]

\RequirePackage{expl3}
\RequirePackage{xparse}
\RequirePackage{l3keys2e}
\RequirePackage{amsmath}
\RequirePackage{amssymb}      % \leqslant / \geqslant in rendered <= / >=
\RequirePackage{tikz}
\usetikzlibrary{shapes.symbols}
\usetikzlibrary{arrows.meta}
% Custom Cloud arrow tip used as the open end of inflow / outflow
% pipes.  The silhouette consists of nine concatenated elliptical
% arcs (\pgfpatharc): each arc starts at the current path point,
% which pgf treats as lying on an ellipse at angle <start>, and
% sweeps through to angle <end>.  Varying the start and end angles
% per arc produces zigzagging half-circle puffs -- the
% characteristic cloud outline.  A single factor \puffscale sets the
% rx and ry of every ellipse and thus the size of all puffs at once.
% The chain has a horizontal span of (4+2*sqrt(2))*rx ~ 6.83*rx;
% with puffscale=0.22 that is ~ 1.5*\pgfarrowlength.  The path
% therefore starts at -0.5*\pgfarrowlength so the silhouette covers
% [-0.5L, L] in local arrow coordinates and the apex lands exactly
% at \pgfarrowlength -- where the built-in tips (Latex, Stealth,
% ...) also place their tip.  The line width is set explicitly to
% 0.3pt, independent of the arrow's line width, so a thick arrow
% does not also get a thick cloud outline.
\pgfdeclarearrow{
  name = Cloud,
  parameters = { \the\pgfarrowlength \the\pgfarrowwidth },
  setup code = {
    \pgfarrowssettipend{\pgfarrowlength}
    \pgfarrowssetbackend{-0.5\pgfarrowlength}
    \pgfarrowssetlineend{-0.5\pgfarrowlength}
    \pgfarrowssetvisualbackend{-0.5\pgfarrowlength}
    \pgfarrowssetvisualtipend{\pgfarrowlength}
  },
  defaults = { length = 1em, width = 1em },
  drawing code = {
    \def\puffscale{0.22}
    \edef\puffrx{\puffscale\pgfarrowlength}
    \edef\puffry{\puffscale\pgfarrowwidth}
    \pgfsetlinewidth{0.3pt}
    \pgfpathmoveto{\pgfqpoint{-0.5\pgfarrowlength}{0pt}}
    \pgfpatharc{180}{ 90}{\puffrx\space and \puffry}
    \pgfpatharc{225}{ 45}{\puffrx\space and \puffry}
    \pgfpatharc{180}{  0}{\puffrx\space and \puffry}
    \pgfpatharc{135}{-45}{\puffrx\space and \puffry}
    \pgfpatharc{ 90}{-90}{\puffrx\space and \puffry}
    \pgfpatharc{ 45}{-135}{\puffrx\space and \puffry}
    \pgfpatharc{  0}{-180}{\puffrx\space and \puffry}
    \pgfpatharc{-45}{-225}{\puffrx\space and \puffry}
    \pgfpatharc{-90}{-180}{\puffrx\space and \puffry}
    \pgfpathclose
    \pgfusepathqfillstroke
  },
}
\RequirePackage{luacode}
\RequirePackage{siunitx}
\RequirePackage{float}        % \diagrammodel uses [H] placement
\RequirePackage{tabularray}   % \textmodel uses longtblr (page-breakable,
                              % works inside floats/subfigure)
% Define a numodel theme for longtblr in \textmodel.  The default
% `normal' middlehead/lasthead prefixes the continuation marker with
% `Table N:' from \tablename/\thetable; we don't issue \caption, so we
% want only the conthead text (which is itself localised via
% \tblrcontheadname, updated by \__numodel_refresh_kw:).  Likewise for
% the foot.  Activated via `theme = {numodel}' in the longtblr options.
\DefTblrTemplate {firsthead}            {numodel} {}
\DefTblrTemplate {middlehead, lasthead} {numodel} { \UseTblrTemplate {conthead} {default} }
\DefTblrTemplate {firstfoot, middlefoot}{numodel} { \UseTblrTemplate {contfoot} {default} }
\DefTblrTemplate {lastfoot}             {numodel} {}
\NewTblrTheme {numodel}
  {
    \SetTblrTemplate {firsthead}            {numodel}
    \SetTblrTemplate {middlehead, lasthead} {numodel}
    \SetTblrTemplate {firstfoot, middlefoot}{numodel}
    \SetTblrTemplate {lastfoot}             {numodel}
  }
% Tighten the default row separation of all tabularray tables: the
% \textmodel rule listing is denser and more compact than the
% tabularray default.  Users who want the default spacing back in
% their own tables can simply issue \SetTblrInner{rowsep=<value>} in
% their document.
\SetTblrInner{rowsep=0pt}
\RequirePackage{numodel-plot}

% ====================================================================
% Non-expl3 helpers: `\nm@def@num@cs` / `\nm@def@qty@cs` / `\nm@def@pre@cs`
% define `\<fullname><suffix>` macros whose body contains literal
% siunitx options.  Defined OUTSIDE \ExplSyntaxOn so the option-string
% colons (`-3:n`, `-0:0`) and other punctuation keep normal catcodes.
% Storing those bodies inside an expl3-context would tokenise `:` as
% letter, which then breaks siunitx's option parser at expansion time.
% ====================================================================
\newcommand{\NumodelDefNumCs}[3]{% #1 = fullname, #2 = start expr, #3 = sigfigs
  \expandafter\gdef\csname #1num\endcsname{%
    \num[evaluate-expression=true,
         round-mode=figures, round-precision=#3,
         exponent-mode=threshold, exponent-thresholds=-3:#3]
        {#2}\relax
  }%
}
\newcommand{\NumodelDefQtyCs}[4]{% #1 = fullname, #2 = expr, #3 = sigfigs, #4 = unit
  \expandafter\gdef\csname #1qty\endcsname{%
    \qty[evaluate-expression=true,
         round-mode=figures, round-precision=#3,
         exponent-mode=threshold, exponent-thresholds=-3:#3]
        {#2}{#4}\relax
  }%
}
\newcommand{\NumodelDefPreCs}[4]{% same parameters as above
  \expandafter\gdef\csname #1pre\endcsname{%
    \qty[evaluate-expression=true,
         round-mode=figures, round-precision=#3,
         prefix-mode=combine-exponent,
         exponent-mode=engineering,
         exponent-thresholds=-0:0]
        {#2}{#4}\relax
  }%
}

\ExplSyntaxOn

% ====================================================================
% Lua module for data storage (O(1) append, O(N) total)
% ====================================================================
% All variable values are stored per step in Lua tables.  This
% replaces the O(N^2) \xdef accumulation for coordinates and makes
% min/max tracking efficient.  After \computemodel the results are
% pushed back to TeX macros.

\begin{luacode*}
local f = kpse.find_file("numodel.lua", "tex")
if f and f ~= "" then
  dofile(f)
else
  tex.error("numodel: cannot find companion Lua module numodel.lua."
         .. " Install it in a directory searched by kpse"
         .. " (e.g. TEXMFHOME/tex/lualatex/numodel/).")
end
\end{luacode*}

% ====================================================================
% Internal data structures
% ====================================================================

\seq_new:N  \g_mvar_names_seq      % all model variable names
\seq_new:N  \g_mvar_start_seq      % names with an initial value (for the table)
\seq_new:N  \g_mrule_seq           % model rules (display strings)
\seq_new:N  \g_mrule_type_seq      % type per display row: rule | cont
\seq_new:N  \g_mrule_calc_seq      % model rules for execution: {name}{expr}
\int_new:N  \g_mrule_counter_int   % rule counter (display numbering)
\tl_new:N   \g_numodel_stop_expr_tl  % stop-condition expression for execution
\int_new:N  \g_numodel_steps_int     % number of executed steps
\int_new:N  \g_numodel_maxiter_int   % safety limit (init via \numodelsetup)
\int_new:N  \g_numodel_gridmaxx_int  % \graphicmodel wrap threshold (0 = off)

% --- Graphic-model infrastructure (Phase 1) ---
\tl_new:N   \l__numodel_tmp_tl       % temporary helper
\tl_new:N   \l__numodel_scratch_tl   % scratch for type/text checks
\tl_new:N   \l__numodel_scratch_y_tl % scratch y-coord (svgy lookup)

% Diagram style for \graphicmodel (see \numodelsetup):
%   tight     -- valve takes the label of the direct inflow variable;
%                the aux/const is moved to the valve position and not
%                drawn as a separate node (compact, default).
%   forrester -- valve has no label; the aux/const stays at its own
%                gridy position; causal arrow from aux/const to the
%                valve.
%   edu       -- combination: valve takes the label *and* a separate
%                aux/const node remains with a causal arrow to it.
\tl_new:N \g__numodel_diagram_style_tl

% Sub-keys that fine-tune the appearance of the Forrester diagram
% (see \numodelsetup).  Empty value = "follow diagram-style default".
%   flowarrow-style    : hollow | filled
%                        forrester -> hollow, tight|edu -> filled
%   valve-style        : valve | circle | edu
%                        forrester -> valve, tight|edu -> edu
%   flowarrow-cloud-tip: true | false
%                        forrester -> true, tight|edu -> false
% Per-variable override for flowarrow-cloud-tip is stored in
% \<name>flowcloud (set via the [flowarrow-cloud-tip=...] key on
% \mvar; empty means "follow global key").
\tl_new:N \g__numodel_flowarrow_style_tl
\tl_new:N \g__numodel_valve_style_tl
\tl_new:N \g__numodel_flowcloud_tl

% Units in the \textmodel initial-values column (see \numodelsetup):
%   true  (default) -- initial value rendered via \<name>qty
%                      (number + unit)
%   false           -- initial value rendered via \<name>num
%                      (number only)
% Per-table override via \textmodel[units=true|false].
\bool_new:N \g__numodel_units_bool

% tabularray environment used by \textmodel (see \numodelsetup{tblrenv}):
%   longtblr -- default; page-breakable; suitable for body text.
%   tblr     -- basic; no page breaks; choose this when \textmodel sits
%               inside an outer environment that already suppresses page
%               breaks (subfigure, minipage, ...).
%   talltblr -- like tblr (no page breaks) but supports captions/notes.
% Per-table override via \textmodel[tblrenv=...].
\tl_new:N \g__numodel_tblrenv_tl

% Decimal separator for \textmodel initial values and \diagrammodel
% tick values (see \numodelsetup{decimal-separator}):
%   point   -- '.' (siunitx output-decimal-marker={.})
%   comma   -- ',' (siunitx output-decimal-marker={,})
% The default follows syntax: english -> point, coachtaal -> comma.
% An explicit decimal-separator key sets
% \g__numodel_dsep_explicit_bool so that a later syntax change no
% longer overrules the choice.
\tl_new:N   \g__numodel_dsep_tl
\bool_new:N \g__numodel_dsep_explicit_bool

% ====================================================================
% Syntax lookup
% ====================================================================
% Determines the language of the text-rendered model.  Affects display
% only (\textmodel, \mruletext, \mstop); \computemodel always uses
% \fpeval internally and is language-agnostic.  Each value of
% \g__numodel_syntax_tl is a tag that selects both the backing
% \__numodel_kw_<tag>_<key>: macros and the file numodel-<tag>.def
% that defines them; the package ships EN (English/XMILE) and NL
% (Dutch/CoachTaal).
%
% The CTAN default is EN.  The initial value is set below via
% \numodelsetup.

\tl_new:N \g__numodel_syntax_tl

% Lookup (fully expandable): \__numodel_kw:n {<key>} returns the
% translation in the current syntax language.  The key also controls
% the surrounding spaces (see keys 'then' vs 'then_nl' etc.).  The
% backing macros \__numodel_kw_<LANG>_<key>: are provided by the
% language file numodel-<LANG>.def loaded by \__numodel_load_syntax_def:n
% (see below).
\cs_new:Npn \__numodel_kw:n #1
  { \use:c { __numodel_kw_ \g__numodel_syntax_tl _ #1 : } }

% Pre-computed translations in tl variables so they are usable via
% \u{} in l3regex replacement text.  (l3regex treats {...} in the
% replacement as brace groups, so a bare \__numodel_kw:n{key} call
% fails there.)
\tl_new:N \g__numodel_kw_if_tl
\tl_new:N \g__numodel_kw_then_tl
\tl_new:N \g__numodel_kw_then_nl_tl
\tl_new:N \g__numodel_kw_else_tl
\tl_new:N \g__numodel_kw_else_nl_tl
\tl_new:N \g__numodel_kw_endif_tl
\tl_new:N \g__numodel_kw_endif_nl_tl
\tl_new:N \g__numodel_kw_and_tl
\tl_new:N \g__numodel_kw_or_tl
\tl_new:N \g__numodel_kw_not_tl
\tl_new:N \g__numodel_kw_sign_tl
\tl_new:N \g__numodel_kw_abs_tl
\tl_new:N \g__numodel_kw_sqrt_tl
\tl_new:N \g__numodel_kw_exp_tl
\tl_new:N \g__numodel_kw_ln_tl
\tl_new:N \g__numodel_kw_sin_tl
\tl_new:N \g__numodel_kw_cos_tl
\tl_new:N \g__numodel_kw_tan_tl
\tl_new:N \g__numodel_kw_asin_tl
\tl_new:N \g__numodel_kw_acos_tl
\tl_new:N \g__numodel_kw_atan_tl
\tl_new:N \g__numodel_kw_min_tl
\tl_new:N \g__numodel_kw_max_tl
\tl_new:N \g__numodel_kw_stop_tl

% Recomputes all tl caches from \g__numodel_syntax_tl.  Must be
% called after any change to the syntax language.
\cs_new_protected:Npn \__numodel_refresh_kw:
  {
    \tl_gset:Ne \g__numodel_kw_if_tl       { \__numodel_kw:n {if}       }
    \tl_gset:Ne \g__numodel_kw_then_tl     { \__numodel_kw:n {then}     }
    \tl_gset:Ne \g__numodel_kw_then_nl_tl  { \__numodel_kw:n {then_nl}  }
    \tl_gset:Ne \g__numodel_kw_else_tl     { \__numodel_kw:n {else}     }
    \tl_gset:Ne \g__numodel_kw_else_nl_tl  { \__numodel_kw:n {else_nl}  }
    \tl_gset:Ne \g__numodel_kw_endif_tl    { \__numodel_kw:n {endif}    }
    \tl_gset:Ne \g__numodel_kw_endif_nl_tl { \__numodel_kw:n {endif_nl} }
    \tl_gset:Ne \g__numodel_kw_and_tl      { \__numodel_kw:n {and}      }
    \tl_gset:Ne \g__numodel_kw_or_tl       { \__numodel_kw:n {or}       }
    \tl_gset:Ne \g__numodel_kw_not_tl      { \__numodel_kw:n {not}      }
    \tl_gset:Ne \g__numodel_kw_sign_tl     { \__numodel_kw:n {sign}     }
    \tl_gset:Ne \g__numodel_kw_abs_tl      { \__numodel_kw:n {abs}      }
    \tl_gset:Ne \g__numodel_kw_sqrt_tl     { \__numodel_kw:n {sqrt}     }
    \tl_gset:Ne \g__numodel_kw_exp_tl      { \__numodel_kw:n {exp}      }
    \tl_gset:Ne \g__numodel_kw_ln_tl       { \__numodel_kw:n {ln}       }
    \tl_gset:Ne \g__numodel_kw_sin_tl      { \__numodel_kw:n {sin}      }
    \tl_gset:Ne \g__numodel_kw_cos_tl      { \__numodel_kw:n {cos}      }
    \tl_gset:Ne \g__numodel_kw_tan_tl      { \__numodel_kw:n {tan}      }
    \tl_gset:Ne \g__numodel_kw_asin_tl     { \__numodel_kw:n {asin}     }
    \tl_gset:Ne \g__numodel_kw_acos_tl     { \__numodel_kw:n {acos}     }
    \tl_gset:Ne \g__numodel_kw_atan_tl     { \__numodel_kw:n {atan}     }
    \tl_gset:Ne \g__numodel_kw_min_tl      { \__numodel_kw:n {min}      }
    \tl_gset:Ne \g__numodel_kw_max_tl      { \__numodel_kw:n {max}      }
    \tl_gset:Ne \g__numodel_kw_stop_tl     { \__numodel_kw:n {stop}     }
    % tabularray continuation markers (used by longtblr in \textmodel)
    \tl_gset:Ne \tblrcontfootname { \__numodel_kw:n {contfoot} }
    \tl_gset:Ne \tblrcontheadname { \__numodel_kw:n {conthead} }
  }

% Helper: \__numodel_kwt:n {<key>} expands to \text{<kw-value>}.
% Meant for use inside \tl_set:Ne / \seq_gput_right:Ne -- the kw is
% inserted during e-expansion while \text remains protected.
\cs_new:Npn \__numodel_kwt:n #1 { \text { \__numodel_kw:n {#1} } }

% Language-file loader.  Each tag <LANG> is backed by a file
% numodel-<LANG>.def installed in a kpse-searched directory; the
% package ships numodel-EN.def and numodel-NL.def, users can drop
% additional files in TEXMFHOME/tex/latex/numodel/ without rebuilding
% the package.
%
% A small alias property maps the historical long names to the
% canonical two-letter tag used in the file name and the internal
% macro names.  Users who add their own file just pass the file's tag
% to syntax=<tag> directly; no alias is required.
\prop_new:N \g__numodel_syntax_aliases_prop
\prop_gput:Nnn \g__numodel_syntax_aliases_prop { english }   { EN }
\prop_gput:Nnn \g__numodel_syntax_aliases_prop { coachtaal } { NL }
\prop_gput:Nnn \g__numodel_syntax_aliases_prop { dutch }     { NL }

% Tags whose .def file has already been input during this run, so the
% lookup only triggers \InputIfFileExists the first time.
\seq_new:N \g__numodel_loaded_syntax_seq

\msg_new:nnn { numodel } { unknown-syntax }
  {
    Cannot~load~syntax~'#1':~file~'numodel-#1.def'~not~found~by~
    kpse.~The~package~ships~with~EN~(English)~and~NL~(Dutch~
    CoachTaal);~drop~your~own~numodel-<LANG>.def~in~
    TEXMFHOME/tex/latex/numodel/~to~add~more~languages.
  }

% \__numodel_load_syntax_def:n {<tag>}
% Inputs numodel-<tag>.def the first time it is requested.  Raises a
% LaTeX error if the file is not found by kpse.
\cs_new_protected:Npn \__numodel_load_syntax_def:n #1
  {
    \seq_if_in:NnF \g__numodel_loaded_syntax_seq {#1}
      {
        \InputIfFileExists { numodel-#1.def }
          { \seq_gput_right:Nn \g__numodel_loaded_syntax_seq {#1} }
          { \msg_error:nnn { numodel } { unknown-syntax } {#1} }
      }
  }

\tl_new:N \l__numodel_syntax_tag_tl

% \__numodel_set_syntax:n {<value>}
% Public-facing setter behind the syntax key.  Resolves aliases,
% loads the matching .def file (once), publishes the canonical tag
% in \g__numodel_syntax_tl, refreshes the keyword cache, and -- when
% the user has not pinned a decimal separator explicitly -- adopts the
% language's preferred decimal mark via the optional
% \__numodel_kw_<LANG>_dsep_default: hook.
\cs_new_protected:Npn \__numodel_set_syntax:n #1
  {
    \prop_get:NnNF \g__numodel_syntax_aliases_prop {#1}
      \l__numodel_syntax_tag_tl
      { \tl_set:Nn \l__numodel_syntax_tag_tl {#1} }
    \exp_args:NV \__numodel_load_syntax_def:n \l__numodel_syntax_tag_tl
    \tl_gset_eq:NN \g__numodel_syntax_tl \l__numodel_syntax_tag_tl
    \__numodel_refresh_kw:
    \bool_if:NF \g__numodel_dsep_explicit_bool
      {
        \cs_if_exist:cTF
          { __numodel_kw_ \g__numodel_syntax_tl _dsep_default: }
          { \tl_gset:Ne \g__numodel_dsep_tl { \__numodel_kw:n {dsep_default} } }
          { \tl_gset:Nn \g__numodel_dsep_tl { point } }
      }
  }

% Applies the decimal separator (only for the duration of the
% surrounding TeX group, so a document-wide \sisetup is left
% untouched).  Calls both siunitx (for \num/\qty in \textmodel
% initial values) and pgfplots' tick styles (for \diagrammodel tick
% labels) so the choice is consistent everywhere.
\cs_new_protected:Npn \__numodel_apply_dsep:
  {
    \str_if_eq:VnTF \g__numodel_dsep_tl { comma }
      {
        \sisetup{ output-decimal-marker = {,} }
        % Append behind the existing numodel/axis style so our
        % xticklabel-style replaces what numodel-plot hard-codes.
        % The \pgfplotsset mutation is \def-based and hence
        % group-local.
        \pgfplotsset
          {
            numodel/axis/.append~style=
              {
                xticklabel~style=
                  { /pgf/number~format/.cd, use~comma, /tikz/.cd } ,
                yticklabel~style=
                  { /pgf/number~format/.cd, use~comma, /tikz/.cd } ,
              }
          }
      }
      {
        \sisetup{ output-decimal-marker = {.} }
        \pgfplotsset
          {
            numodel/axis/.append~style=
              {
                xticklabel~style=
                  { /pgf/number~format/.cd, use~period, /tikz/.cd } ,
                yticklabel~style=
                  { /pgf/number~format/.cd, use~period, /tikz/.cd } ,
              }
          }
      }
  }

% ====================================================================
% Configuration command \numodelsetup
% ====================================================================
% Runtime API for settings.  Keys:
%   syntax     -- language tag (file numodel-<tag>.def must be on kpse
%                 path).  Built-in: EN, NL.  Legacy aliases: english,
%                 coachtaal, dutch.
%   maxiter    -- safety limit for \computemodel (default 20000)
%   graphscalex -- horizontal grid spacing \dgridx for Forrester
%                  diagrams (default 2)
%   graphscaley -- vertical grid spacing \dgridy for Forrester
%                  diagrams (default 2)

% Helper for choice-with-empty-reset: validates #4 against #3 (a
% comma-separated whitelist) and writes it into #1 (a tl).  An empty
% value clears the tl (which means "follow the diagram-style
% default" for the flowarrow-style / valve-style /
% flowarrow-cloud-tip keys).  Any other value triggers a warning.
\msg_new:nnn { numodel } { bad-choice }
  { Key~'#1'~accepts~only~#2,~or~empty~to~reset.~Got:~'#3'. }

\cs_new_protected:Npn \__numodel_setup_choice:Nnnn #1 #2 #3 #4
  {
    \tl_if_blank:nTF {#4}
      { \tl_gclear:N #1 }
      {
        \clist_if_in:nnTF {#3} {#4}
          { \tl_gset:Nn #1 {#4} }
          { \msg_warning:nnnnn { numodel } { bad-choice } {#2} {#3} {#4} }
      }
  }

% Local-tl variant: same validation, local set/clear.  Used by the
% \graphicmodel one-render override and the per-\mvar override.
\cs_new_protected:Npn \__numodel_local_choice:Nnnn #1 #2 #3 #4
  {
    \tl_if_blank:nTF {#4}
      { \tl_clear:N #1 }
      {
        \clist_if_in:nnTF {#3} {#4}
          { \tl_set:Nn #1 {#4} }
          { \msg_warning:nnnnn { numodel } { bad-choice } {#2} {#3} {#4} }
      }
  }

\keys_define:nn { numodel / setup }
  {
    syntax              .code:n =
      { \__numodel_set_syntax:n {#1} },
    maxiter             .int_gset:N = \g_numodel_maxiter_int,
    graphscalex         .code:n = { \tl_gset:Nn \dgridx {#1} },
    graphscaley         .code:n = { \tl_gset:Nn \dgridy {#1} },
    gridmaxx            .int_gset:N = \g_numodel_gridmaxx_int,
    diagram-style              .choice:,
    diagram-style / tight      .code:n =
      { \tl_gset:Nn \g__numodel_diagram_style_tl { tight } },
    diagram-style / forrester  .code:n =
      { \tl_gset:Nn \g__numodel_diagram_style_tl { forrester } },
    diagram-style / edu        .code:n =
      { \tl_gset:Nn \g__numodel_diagram_style_tl { edu } },
    flowarrow-style     .code:n =
      { \__numodel_setup_choice:Nnnn \g__numodel_flowarrow_style_tl
          { flowarrow-style } { hollow , filled } {#1} },
    valve-style         .code:n =
      { \__numodel_setup_choice:Nnnn \g__numodel_valve_style_tl
          { valve-style } { valve , circle , edu } {#1} },
    flowarrow-cloud-tip .code:n =
      { \__numodel_setup_choice:Nnnn \g__numodel_flowcloud_tl
          { flowarrow-cloud-tip } { true , false } {#1} },
    units                      .choice:,
    units / true               .code:n =
      { \bool_gset_true:N  \g__numodel_units_bool },
    units / false              .code:n =
      { \bool_gset_false:N \g__numodel_units_bool },
    tblrenv                    .choice:,
    tblrenv / tblr             .code:n =
      { \tl_gset:Nn \g__numodel_tblrenv_tl { tblr } },
    tblrenv / longtblr         .code:n =
      { \tl_gset:Nn \g__numodel_tblrenv_tl { longtblr } },
    tblrenv / talltblr         .code:n =
      { \tl_gset:Nn \g__numodel_tblrenv_tl { talltblr } },
    decimal-separator          .choice:,
    decimal-separator / comma  .code:n =
      {
        \tl_gset:Nn \g__numodel_dsep_tl { comma }
        \bool_gset_true:N \g__numodel_dsep_explicit_bool
      },
    decimal-separator / point  .code:n =
      {
        \tl_gset:Nn \g__numodel_dsep_tl { point }
        \bool_gset_true:N \g__numodel_dsep_explicit_bool
      },
  }

\NewDocumentCommand{\numodelsetup}{ m }
  { \keys_set:nn { numodel / setup } {#1} }

% Defaults (this call also initialises \g__numodel_syntax_tl,
% \g_numodel_maxiter_int, \dgridx, \dgridy and
% \g__numodel_diagram_style_tl).
\numodelsetup
  {
    syntax        = EN,
    maxiter       = 20000,
    graphscalex   = 2,
    graphscaley   = 2,
    diagram-style = tight,
    units         = true,
    tblrenv       = longtblr,
  }

% Package-time options.  \usepackage[syntax=NL]{numodel} delegates
% to the same key infrastructure as \numodelsetup.
\keys_define:nn { numodel / pkg }
  {
    syntax              .meta:nn = { numodel / setup }{ syntax              = #1 },
    maxiter             .meta:nn = { numodel / setup }{ maxiter             = #1 },
    graphscalex         .meta:nn = { numodel / setup }{ graphscalex         = #1 },
    graphscaley         .meta:nn = { numodel / setup }{ graphscaley         = #1 },
    gridmaxx            .meta:nn = { numodel / setup }{ gridmaxx            = #1 },
    diagram-style       .meta:nn = { numodel / setup }{ diagram-style       = #1 },
    flowarrow-style     .meta:nn = { numodel / setup }{ flowarrow-style     = #1 },
    valve-style         .meta:nn = { numodel / setup }{ valve-style         = #1 },
    flowarrow-cloud-tip .meta:nn = { numodel / setup }{ flowarrow-cloud-tip = #1 },
    units               .meta:nn = { numodel / setup }{ units               = #1 },
    tblrenv             .meta:nn = { numodel / setup }{ tblrenv             = #1 },
    decimal-separator   .meta:nn = { numodel / setup }{ decimal-separator   = #1 },
  }
\ProcessKeyOptions [ numodel / pkg ]

% ====================================================================
% Prefix system
% ====================================================================
% Each model lives under a prefix (a short string).  The "live state"
% sits in the \g_mvar_*, \g_mrule_*, \g_numodel_* variables above;
% \newmodelprefix and \switchmodelprefix swap that state to/from
% per-prefix backup storage.
%
% - \g_numodel_current_prefix_tl: current prefix (empty at package load)
% - \g_numodel_prefixes_seq: list of all registered prefixes
% - \l__numodel_cmd_prefix_tl: prefix passed via [prefix=...] key
% - \l__numodel_eff_prefix_tl: effective prefix for the current command

\tl_new:N  \g_numodel_current_prefix_tl
\seq_new:N \g_numodel_prefixes_seq
\tl_new:N  \l__numodel_cmd_prefix_tl
\tl_new:N  \l__numodel_eff_prefix_tl
\tl_new:N  \l__numodel_fullname_tl

% Per-prefix storage: on swap the current state is copied to
% \g__numodel_<P>_<slot>_{seq,tl,int,prop}.  Load goes the other way.
\cs_new_protected:Npn \__numodel_save_state:n #1
  {
    \seq_gset_eq:cN { g__numodel_ #1 _vars_seq      } \g_mvar_names_seq
    \seq_gset_eq:cN { g__numodel_ #1 _starts_seq    } \g_mvar_start_seq
    \seq_gset_eq:cN { g__numodel_ #1 _rules_seq     } \g_mrule_seq
    \seq_gset_eq:cN { g__numodel_ #1 _ruletypes_seq } \g_mrule_type_seq
    \seq_gset_eq:cN { g__numodel_ #1 _rulecalc_seq  } \g_mrule_calc_seq
    \int_gset:cn    { g__numodel_ #1 _rulecounter_int }
      { \int_use:N \g_mrule_counter_int }
    \tl_gset_eq:cN  { g__numodel_ #1 _stopexpr_tl   } \g_numodel_stop_expr_tl
    \int_gset:cn    { g__numodel_ #1 _steps_int     }
      { \int_use:N \g_numodel_steps_int }
  }

\cs_new_protected:Npn \__numodel_load_state:n #1
  {
    \seq_gset_eq:Nc \g_mvar_names_seq  { g__numodel_ #1 _vars_seq      }
    \seq_gset_eq:Nc \g_mvar_start_seq  { g__numodel_ #1 _starts_seq    }
    \seq_gset_eq:Nc \g_mrule_seq       { g__numodel_ #1 _rules_seq     }
    \seq_gset_eq:Nc \g_mrule_type_seq  { g__numodel_ #1 _ruletypes_seq }
    \seq_gset_eq:Nc \g_mrule_calc_seq  { g__numodel_ #1 _rulecalc_seq  }
    \int_gset:Nn \g_mrule_counter_int
      { \int_use:c { g__numodel_ #1 _rulecounter_int } }
    \tl_gset_eq:Nc \g_numodel_stop_expr_tl  { g__numodel_ #1 _stopexpr_tl }
    \int_gset:Nn \g_numodel_steps_int
      { \int_use:c { g__numodel_ #1 _steps_int } }
  }

\cs_new_protected:Npn \__numodel_clear_state:
  {
    \seq_gclear:N  \g_mvar_names_seq
    \seq_gclear:N  \g_mvar_start_seq
    \seq_gclear:N  \g_mrule_seq
    \seq_gclear:N  \g_mrule_type_seq
    \seq_gclear:N  \g_mrule_calc_seq
    \int_gzero:N   \g_mrule_counter_int
    \tl_gclear:N   \g_numodel_stop_expr_tl
    \int_gzero:N   \g_numodel_steps_int
  }

% Initialise per-prefix backup variables (once, at \newmodelprefix).
\cs_new_protected:Npn \__numodel_init_backup:n #1
  {
    \seq_new:c  { g__numodel_ #1 _vars_seq      }
    \seq_new:c  { g__numodel_ #1 _starts_seq    }
    \seq_new:c  { g__numodel_ #1 _rules_seq     }
    \seq_new:c  { g__numodel_ #1 _ruletypes_seq }
    \seq_new:c  { g__numodel_ #1 _rulecalc_seq  }
    \int_new:c  { g__numodel_ #1 _rulecounter_int }
    \tl_new:c   { g__numodel_ #1 _stopexpr_tl   }
    \int_new:c  { g__numodel_ #1 _steps_int     }
  }

% Public commands for prefix management
\NewDocumentCommand{\newmodelprefix}{ m }
  {
    \typeout{NEWPREFIX~start:~#1}
    \seq_if_in:NnTF \g_numodel_prefixes_seq {#1}
      { \msg_error:nnn { numodel } { prefix-exists } {#1} }
      {
        \typeout{NEWPREFIX~save~current:~\g_numodel_current_prefix_tl}
        % Save current state under the previous prefix (if any)
        \tl_if_empty:NF \g_numodel_current_prefix_tl
          { \exp_args:NV \__numodel_save_state:n \g_numodel_current_prefix_tl }
        \typeout{NEWPREFIX~register}
        % Register new prefix + init backup vars + Lua state
        \seq_gput_right:Nn \g_numodel_prefixes_seq {#1}
        \typeout{NEWPREFIX~init_backup}
        \__numodel_init_backup:n {#1}
        \typeout{NEWPREFIX~lua_init}
        \directlua{ numodel.init_prefix("#1") }
        \typeout{NEWPREFIX~clear_state}
        % Clear current state and set the prefix
        \__numodel_clear_state:
        \typeout{NEWPREFIX~set_current}
        \tl_gset:Nn \g_numodel_current_prefix_tl {#1}
        \typeout{NEWPREFIX~done:~#1}
      }
  }

\NewDocumentCommand{\switchmodelprefix}{ m }
  {
    \seq_if_in:NnTF \g_numodel_prefixes_seq {#1}
      {
        % Save current state (if any), load the new one
        \tl_if_empty:NF \g_numodel_current_prefix_tl
          { \exp_args:NV \__numodel_save_state:n \g_numodel_current_prefix_tl }
        \__numodel_load_state:n {#1}
        \tl_gset:Nn \g_numodel_current_prefix_tl {#1}
      }
      { \msg_error:nnn { numodel } { prefix-unknown } {#1} }
  }

% Iterate over *all* variables from *all* prefixes.
% #1 = code that uses ##1 = full variable name (prefix + short).
% Public registration API for external tools such as worksheet.tex.
\NewDocumentCommand{\NumodelForEachVar}{ +m }
  {
    \seq_map_inline:Nn \g_numodel_prefixes_seq
      {
        \seq_map_inline:cn { g__numodel_ ##1 _vars_seq } {#1}
      }
  }

% Resolve the effective prefix for a command call.  Takes
% [prefix=<p>] from the keyval argument; otherwise falls back to the
% current prefix.  Stores the result in \l__numodel_eff_prefix_tl.
\keys_define:nn { numodel / cmd }
  {
    prefix        .tl_set:N = \l__numodel_cmd_prefix_tl,
    % \graphicmodel-only: temporary override of the global
    % diagram-style.  Cleared after every \graphicmodel call.
    diagram-style              .choice:,
    diagram-style / tight      .code:n =
      { \tl_set:Nn \l__numodel_cmd_diagstyle_tl { tight } },
    diagram-style / forrester  .code:n =
      { \tl_set:Nn \l__numodel_cmd_diagstyle_tl { forrester } },
    diagram-style / edu        .code:n =
      { \tl_set:Nn \l__numodel_cmd_diagstyle_tl { edu } },
    % \graphicmodel-only: temporary overrides of the corresponding
    % global flowarrow/valve/cloud keys.  Empty after every call.
    flowarrow-style     .code:n =
      { \__numodel_local_choice:Nnnn \l__numodel_cmd_flowarrow_tl
          { flowarrow-style } { hollow , filled } {#1} },
    valve-style         .code:n =
      { \__numodel_local_choice:Nnnn \l__numodel_cmd_valve_tl
          { valve-style } { valve , circle , edu } {#1} },
    flowarrow-cloud-tip .code:n =
      { \__numodel_local_choice:Nnnn \l__numodel_cmd_flowcloud_tl
          { flowarrow-cloud-tip } { true , false } {#1} },
    % \textmodel-only: temporary override of the global units bool.
    % Cleared after every \textmodel call.
    units                      .choice:,
    units / true               .code:n =
      { \tl_set:Nn \l__numodel_cmd_units_tl { true  } },
    units / false              .code:n =
      { \tl_set:Nn \l__numodel_cmd_units_tl { false } },
    % \textmodel-only: temporary override of the global tblrenv tl.
    % Cleared after every \textmodel call.
    tblrenv                    .choice:,
    tblrenv / tblr             .code:n =
      { \tl_set:Nn \l__numodel_cmd_tblrenv_tl { tblr } },
    tblrenv / longtblr         .code:n =
      { \tl_set:Nn \l__numodel_cmd_tblrenv_tl { longtblr } },
    tblrenv / talltblr         .code:n =
      { \tl_set:Nn \l__numodel_cmd_tblrenv_tl { talltblr } },
  }
\tl_new:N \l__numodel_cmd_diagstyle_tl
\tl_new:N \l__numodel_cmd_flowarrow_tl
\tl_new:N \l__numodel_cmd_valve_tl
\tl_new:N \l__numodel_cmd_flowcloud_tl
\tl_new:N \l__numodel_cmd_units_tl
\tl_new:N \l__numodel_cmd_tblrenv_tl

\cs_new_protected:Npn \__numodel_resolve_prefix:n #1
  {
    \tl_clear:N \l__numodel_cmd_prefix_tl
    \keys_set:nn { numodel / cmd } {#1}
    \tl_if_empty:NTF \l__numodel_cmd_prefix_tl
      { \tl_set_eq:NN \l__numodel_eff_prefix_tl \g_numodel_current_prefix_tl }
      { \tl_set_eq:NN \l__numodel_eff_prefix_tl \l__numodel_cmd_prefix_tl }
  }

% Execute code under a temporarily different prefix (via state swap).
% #1 = target prefix (tl, in variable)
% #2 = code
\cs_new_protected:Npn \__numodel_with_prefix:Nn #1 #2
  {
    \tl_if_eq:NNTF #1 \g_numodel_current_prefix_tl
      { #2 }    % already the current prefix; no swap needed
      {
        \tl_set_eq:NN \l__numodel_saved_prefix_tl \g_numodel_current_prefix_tl
        \exp_args:NV \switchmodelprefix #1
        #2
        \exp_args:NV \switchmodelprefix \l__numodel_saved_prefix_tl
      }
  }
\tl_new:N \l__numodel_saved_prefix_tl

% Build the full name = <prefix><shortname>.  Stored in
% \l__numodel_fullname_tl.
\cs_new_protected:Npn \__numodel_set_fullname:n #1
  {
    \tl_set:Ne \l__numodel_fullname_tl { \l__numodel_eff_prefix_tl #1 }
  }

% Keys for the optional argument of \mvar (grid position +
% initial-value aliases + prefix).
\tl_new:N \l__numodel_alias_tl
\tl_new:N \l__numodel_aliasleft_tl
\tl_new:N \l__numodel_aliasright_tl

\keys_define:nn { numodel / mvar }
  {
    prefix      .tl_set:N = \l__numodel_cmd_prefix_tl ,
    gridx       .tl_set:N = \l__numodel_gridx_tl ,
    gridy       .tl_set:N = \l__numodel_gridy_tl ,
    alias       .tl_set:N = \l__numodel_alias_tl ,
    aliasleft   .tl_set:N = \l__numodel_aliasleft_tl ,
    aliasright  .tl_set:N = \l__numodel_aliasright_tl ,
    flowarrow-cloud-tip .code:n =
      { \__numodel_local_choice:Nnnn \l__numodel_mvar_flowcloud_tl
          { flowarrow-cloud-tip } { true , false } {#1} },
  }
\tl_new:N \l__numodel_mvar_flowcloud_tl

% ====================================================================
% \mvar[keys]{name}{display}{startvalue}{unit}{sig}{type}
% ====================================================================
% Declare a model variable.  Generates the following macros:
%
%   \<name>      -- numeric value (\edef + \fpeval), or warning if empty
%   \<name>text  -- display name for the model table (e.g. F_{res})
%   \<name>unit  -- SI unit (e.g. kN)
%   \<name>unitraw - raw SI unit (e.g. \kilo\N)
%   \<name>sign  -- number of significant figures
%   \<name>type  -- variable type: stock, constant, aux, system
%                   (Dutch synonyms voorraad/constante/hulp/systeem
%                   are normalised to the canonical English form)
%   \<name>coord -- coordinate list (empty; filled by \computemodel)
%   \<name>gridx -- x position in the graphic model (-1 = auto)
%   \<name>gridy -- y position in the graphic model (-1 = auto)
%   \<name>num   -- number with significance (via \num)
%   \<name>qty   -- number + unit (via \qty)
%   \<name>pre   -- engineering-prefix notation (via \qty)
%   \<name>alias       -- replaces the whole initial-value cell
%                         (empty = default)
%   \<name>aliasleft   -- replaces the left symbol in the initial value
%                         (empty = default)
%   \<name>aliasright  -- replaces the right number in the initial value
%                         (empty = default)
%
% Keys accepted in [#1]:
%   gridx, gridy          -- position in the graphic model
%   alias                 -- replace the entire initial-value cell
%                            (in math mode)
%   aliasleft             -- replace just the left symbol
%   aliasright            -- replace just the right value
%
% The base value is registered in \g_defqty_names_seq so that
% \includeimage expands it automatically.
%
% With an empty start value (#3 blank), the base value is not
% numerically defined; using it issues a warning.  Useful for
% auxiliary variables computed by \mrule.
%
% Example:
%   \mvar{modM}{m}{80}{\kg}{2}{constant}
%   \mvar{modFres}{F_{res}}{}{\N}{2}{aux}
%   \mvar[aliasright=\cdots]{modX}{x}{0}{\m}{3}{stock}
%   \mvar[alias={x \text{?}}]{modX}{x}{0}{\m}{3}{stock}

% Normalise the variable type (sixth \mvar argument) to its
% canonical English form.  Accepts the canonical names
% {stock, constant, aux, system} and the Dutch synonyms
% {voorraad, constante, hulp, systeem}.  An unknown value is left
% as-is and a warning is issued.  The result is returned through
% \l__numodel_type_tl.
\tl_new:N \l__numodel_type_tl
\cs_new_protected:Npn \__numodel_normalize_type:n #1
  {
    \str_case:nnF {#1}
      {
        { stock     } { \tl_set:Nn \l__numodel_type_tl { stock    } }
        { constant  } { \tl_set:Nn \l__numodel_type_tl { constant } }
        { aux       } { \tl_set:Nn \l__numodel_type_tl { aux      } }
        { system    } { \tl_set:Nn \l__numodel_type_tl { system   } }
        { voorraad  } { \tl_set:Nn \l__numodel_type_tl { stock    } }
        { constante } { \tl_set:Nn \l__numodel_type_tl { constant } }
        { hulp      } { \tl_set:Nn \l__numodel_type_tl { aux      } }
        { systeem   } { \tl_set:Nn \l__numodel_type_tl { system   } }
      }
      {
        \msg_warning:nne { numodel } { unknown-type } {#1}
        \tl_set:Nn \l__numodel_type_tl {#1}
      }
  }

\makeatletter
\NewDocumentCommand{\mvar}{ O{} m m m m m m }{%
  \typeout{MVAR~start:~#2}
  % Parse optional keys (prefix, gridx, gridy, alias, aliasleft, aliasright,
  % flowarrow-cloud-tip)
  \tl_set:Nn \l__numodel_gridx_tl { -1 }
  \tl_set:Nn \l__numodel_gridy_tl { -1 }
  \tl_clear:N \l__numodel_alias_tl
  \tl_clear:N \l__numodel_aliasleft_tl
  \tl_clear:N \l__numodel_aliasright_tl
  \tl_clear:N \l__numodel_cmd_prefix_tl
  \tl_clear:N \l__numodel_mvar_flowcloud_tl
  \typeout{MVAR~before-keys-set}
  \keys_set:nn { numodel / mvar } {#1}
  \typeout{MVAR~after-keys-set}
  \__numodel_resolve_eff_prefix:
  \typeout{MVAR~eff:~\l__numodel_eff_prefix_tl}
  \__numodel_with_prefix:Nn \l__numodel_eff_prefix_tl
    { \__numodel_mvar_body:nnnnnn {#2}{#3}{#4}{#5}{#6}{#7} }%
  \typeout{MVAR~done:~#2}
}
\makeatother

% Resolve the effective prefix from \l__numodel_cmd_prefix_tl.
\cs_new_protected:Npn \__numodel_resolve_eff_prefix:
  {
    \tl_if_empty:NTF \l__numodel_cmd_prefix_tl
      { \tl_set_eq:NN \l__numodel_eff_prefix_tl \g_numodel_current_prefix_tl }
      { \tl_set_eq:NN \l__numodel_eff_prefix_tl \l__numodel_cmd_prefix_tl }
  }

% The \mvar body runs in the context of the right prefix (live state
% has already been swapped).
\cs_new_protected:Npn \__numodel_mvar_body:nnnnnn #1 #2 #3 #4 #5 #6
  {
    % Guard: current prefix must not be empty (\newmodelprefix required)
    \tl_if_empty:NT \g_numodel_current_prefix_tl
      { \msg_error:nn { numodel } { no-prefix } }
    % Full name = current prefix + short name (#1)
    \tl_set:Ne \l__numodel_fullname_tl { \g_numodel_current_prefix_tl #1 }
    \typeout{MVAR~body~fullname:~\l__numodel_fullname_tl}
    % Register in the per-prefix seq.  The defqty registration is
    % optional: it only fires when the user has loaded the project-
    % specific 'defqty' system (used for worksheet expansion).
    % Standalone, the package works without that seq.
    \seq_gput_right:NV \g_mvar_names_seq   \l__numodel_fullname_tl
    \cs_if_exist:NT \g_defqty_names_seq
      { \seq_gput_right:NV \g_defqty_names_seq \l__numodel_fullname_tl }
    \directlua{ numodel.register(
      "\g_numodel_current_prefix_tl",
      "\l__numodel_fullname_tl") }
    \cs_if_exist:cT { \l__numodel_fullname_tl }
      { \msg_warning:nne { numodel } { redef } { \l__numodel_fullname_tl } }
    \tl_if_blank:nTF {#3}
      { \cs_gset:cpe { \l__numodel_fullname_tl } { \fp_eval:n { 0 } } }
      {
        \cs_gset:cpe { \l__numodel_fullname_tl } { \fp_eval:n {#3} }
        \seq_gput_right:NV \g_mvar_start_seq \l__numodel_fullname_tl
      }
    \cs_gset:cpn { \l__numodel_fullname_tl text }   {#2}
    \cs_gset:cpn { \l__numodel_fullname_tl unit }   { \unit{#4} }
    \cs_gset:cpn { \l__numodel_fullname_tl unitraw } {#4}
    \cs_gset:cpn { \l__numodel_fullname_tl sign }   {#5}
    \__numodel_normalize_type:n {#6}
    \cs_gset:cpe { \l__numodel_fullname_tl type }
      { \tl_use:N \l__numodel_type_tl }
    \cs_gset:cpn { \l__numodel_fullname_tl min }    { inf }
    \cs_gset:cpn { \l__numodel_fullname_tl max }    { -inf }
    \cs_gset:cpe { \l__numodel_fullname_tl gridx }  { \tl_use:N \l__numodel_gridx_tl }
    \cs_gset:cpe { \l__numodel_fullname_tl gridy }  { \tl_use:N \l__numodel_gridy_tl }
    % Preserve the original user input so that \__numodel_build_graphic:
    % can reset auto-placed positions to -1 on a second invocation.
    \cs_gset:cpe { \l__numodel_fullname_tl gridxinit } { \tl_use:N \l__numodel_gridx_tl }
    \cs_gset:cpe { \l__numodel_fullname_tl gridyinit } { \tl_use:N \l__numodel_gridy_tl }
    % Pillar A — Lua-side meta for compute_layout (additive in A1).
    % Detokenize text so that \luaescapestring works on bare chars.
    \tl_set:Ne \l__numodel_scratch_tl { \detokenize {#2} }
    \directlua{ numodel.set_meta(
      "\g_numodel_current_prefix_tl",
      "\l__numodel_fullname_tl",
      { type  = "\tl_use:N \l__numodel_type_tl",
        text  = "\luaescapestring{\l__numodel_scratch_tl}",
        gridx = \tl_use:N \l__numodel_gridx_tl,
        gridy = \tl_use:N \l__numodel_gridy_tl }) }
    \cs_gset:cpe { \l__numodel_fullname_tl alias }      { \exp_not:V \l__numodel_alias_tl }
    \cs_gset:cpe { \l__numodel_fullname_tl aliasleft }  { \exp_not:V \l__numodel_aliasleft_tl }
    \cs_gset:cpe { \l__numodel_fullname_tl aliasright } { \exp_not:V \l__numodel_aliasright_tl }
    \cs_gset:cpe { \l__numodel_fullname_tl flowcloud }  { \exp_not:V \l__numodel_mvar_flowcloud_tl }
    \tl_if_blank:nF {#3}
      {
        \exp_args:NV \NumodelDefNumCs \l__numodel_fullname_tl {\fp_eval:n {#3}} {#5}
        \exp_args:NV \NumodelDefQtyCs \l__numodel_fullname_tl {\fp_eval:n {#3}} {#5} {#4}
        \exp_args:NV \NumodelDefPreCs \l__numodel_fullname_tl {\fp_eval:n {#3}} {#5} {#4}
      }
  }

% ====================================================================
% Warning messages
% ====================================================================

\msg_new:nnn { numodel } { redef }
  { Model~variable~'#1'~is~being~redefined. }
\msg_new:nnn { numodel } { empty-use }
  { Model~variable~'#1'~has~no~start~value~and~is~used~before~assignment. }
\msg_new:nnn { numodel } { maxiter }
  { computemodel~stopped~after~\int_use:N \g_numodel_maxiter_int ~iterations~
    (safety~limit).~Check~stop~condition. }
\msg_new:nnn { numodel } { no-stop }
  { computemodel:~no~stop~condition~defined.~Use~\token_to_str:N \mstop\space
    before~\token_to_str:N \computemodel. }
\msg_new:nnn { numodel } { prefix-exists }
  { Model~prefix~'#1'~is~already~registered.~
    Use~\token_to_str:N \switchmodelprefix\space to~switch~to~an~existing~prefix. }
\msg_new:nnn { numodel } { prefix-unknown }
  { Model~prefix~'#1'~is~not~registered.~
    Use~\token_to_str:N \newmodelprefix\space to~create~it~first. }
\msg_new:nnn { numodel } { no-prefix }
  { No~current~model~prefix.~
    Call~\token_to_str:N \newmodelprefix{<name>}\space before~using~model~commands. }
\msg_new:nnn { numodel } { unknown-type }
  { Unknown~variable~type~'#1'.~
    Use~one~of~stock,~constant,~aux,~system~
    (or~the~Dutch~aliases~voorraad,~constante,~hulp,~systeem). }
\msg_new:nnn { numodel } { unit-mismatch }
  { Variable~'#1'~has~a~different~unit~from~the~first~variable~
    (expected~'#2').~Skipping~it~in~\token_to_str:N \diagrammodel\space
    so~the~remaining~series~share~a~common~y-axis. }

% ====================================================================
% Display helpers
% ====================================================================
% Translate computational expressions into syllabus-style display.
%
% Automatic translations:
%   \modXxx       -> display name (via \<name>text)
%   *             -> \cdot
%   >=  <=        -> \geqslant  \leqslant
%   sign(...)     -> SIGN(...)   / Teken(...) in coachtaal
%   abs(...)      -> ABS(...)    / Abs(...)   in coachtaal
%   sqrt(...)     -> SQRT(...)   / Sqrt(...)   in coachtaal
%   exp(...)      -> EXP(...)    / Exp(...)    in coachtaal
%   ln(...)       -> LN(...)     / Ln(...)     in coachtaal
%   sin(...)      -> SIN(...)    / Sin(...)    in coachtaal
%   cos(...)      -> COS(...)    / Cos(...)    in coachtaal
%   tan(...)      -> TAN(...)    / Tan(...)    in coachtaal
%   asin(...)     -> ARCSIN(...) / Arcsin(...) in coachtaal
%   acos(...)     -> ARCCOS(...) / Arccos(...) in coachtaal
%   &&            -> AND          / EN          in coachtaal
%   ||            -> OR           / OF          in coachtaal
%   cond ? a : b  -> IF cond THEN ... ELSE ... ENDIF (or coachtaal eq.)

\tl_new:N \l__numodel_display_tl
\tl_new:N \l__numodel_lhs_tl
\tl_new:N \l__numodel_rhs_tl
\tl_new:N \l__numodel_cond_tl
\tl_new:N \l__numodel_true_tl
\tl_new:N \l__numodel_false_tl

\cs_new_protected:Npn \__numodel_vars_to_display:N #1
  {
    \typeout{VTD~input:~\tl_to_str:N #1}
    \typeout{VTD~seq:~\seq_use:Nn \g_mvar_names_seq {|}}
    % Variable names -> display names
    \seq_map_inline:Nn \g_mvar_names_seq
      {
        \typeout{VTD~iter:~##1}
        \cs_if_exist:cT { ##1 text }
          {
            \tl_set:Ne \l_tmpa_tl { \use:c { ##1 text } }
            \typeout{VTD~replace~\c{##1}~with~\l_tmpa_tl}
            \regex_replace_all:nnN { \c{##1} } { \u{l_tmpa_tl} } #1
          }
      }
    % Multi-character numeric superscripts.  \fp_eval:n accepts
    % '^1.2' (and '^12', '^-3', ...), but TeX math mode binds only
    % the next single token to ^, so 'T^1.2' would typeset as
    % 'T^{1}.2' and 'T^12' as 'T^{1}2'.  Wrap any numeric exponent
    % literal in braces so display reads 'T^{1.2}', 'T^{12}', etc.
    % Explicitly braced (^{...}) or parenthesised (^(...)) exponents
    % are left untouched.
    \regex_replace_all:nnN { \^ ( [+\-]? [0-9]+ (?: \. [0-9]+)? ) }
      { \^ \cB\{ \1 \cE\} } #1
    % Arithmetic operators
    \regex_replace_all:nnN { \* } { \c{cdot} \x{20} } #1
    \regex_replace_all:nnN { >= } { \c{geqslant} \x{20} } #1
    \regex_replace_all:nnN { <= } { \c{leqslant} \x{20} } #1
    % Functions (syllabus notation).  Order matters: asin/acos must
    % match before sin/cos, otherwise the inner sin/cos gets replaced
    % first and the outer arc- prefix is left dangling.
    \regex_replace_all:nnN { sign \( }
      { \c{text}\cB\{ \u{g__numodel_kw_sign_tl}\cE\}( } #1
    \regex_replace_all:nnN { abs \( }
      { \c{text}\cB\{ \u{g__numodel_kw_abs_tl}\cE\}( } #1
    \regex_replace_all:nnN { sqrt \( }
      { \c{text}\cB\{ \u{g__numodel_kw_sqrt_tl}\cE\}( } #1
    \regex_replace_all:nnN { exp \( }
      { \c{text}\cB\{ \u{g__numodel_kw_exp_tl}\cE\}( } #1
    \regex_replace_all:nnN { ln \( }
      { \c{text}\cB\{ \u{g__numodel_kw_ln_tl}\cE\}( } #1
    \regex_replace_all:nnN { asin \( }
      { \c{text}\cB\{ \u{g__numodel_kw_asin_tl}\cE\}( } #1
    \regex_replace_all:nnN { acos \( }
      { \c{text}\cB\{ \u{g__numodel_kw_acos_tl}\cE\}( } #1
    \regex_replace_all:nnN { atan \( }
      { \c{text}\cB\{ \u{g__numodel_kw_atan_tl}\cE\}( } #1
    \regex_replace_all:nnN { sin \( }
      { \c{text}\cB\{ \u{g__numodel_kw_sin_tl}\cE\}( } #1
    \regex_replace_all:nnN { cos \( }
      { \c{text}\cB\{ \u{g__numodel_kw_cos_tl}\cE\}( } #1
    \regex_replace_all:nnN { \b tan \( }
      { \c{text}\cB\{ \u{g__numodel_kw_tan_tl}\cE\}( } #1
    % Reciprocal trig functions: rendered as 1/<base function>.
    \regex_replace_all:nnN { cot \( }
      { 1/\c{text}\cB\{ \u{g__numodel_kw_tan_tl}\cE\}( } #1
    \regex_replace_all:nnN { csc \( }
      { 1/\c{text}\cB\{ \u{g__numodel_kw_sin_tl}\cE\}( } #1
    \regex_replace_all:nnN { sec \( }
      { 1/\c{text}\cB\{ \u{g__numodel_kw_cos_tl}\cE\}( } #1
    % Two-argument extrema.
    \regex_replace_all:nnN { min \( }
      { \c{text}\cB\{ \u{g__numodel_kw_min_tl}\cE\}( } #1
    \regex_replace_all:nnN { max \( }
      { \c{text}\cB\{ \u{g__numodel_kw_max_tl}\cE\}( } #1
    % Logical operators (syllabus notation)
    \regex_replace_all:nnN { \&\& }
      { \c{text}\cB\{\u{g__numodel_kw_and_tl}\cE\} } #1
    \regex_replace_all:nnN { \|\| }
      { \c{text}\cB\{\u{g__numodel_kw_or_tl}\cE\} } #1
  }

% Ternary detection: cond ? true_expr : false_expr
% Converted to: IF cond THEN lhs = true ELSE lhs = false ENDIF
% (or the coachtaal equivalent).  Supports only flat (non-nested)
% ternaries.
\bool_new:N \l__numodel_is_ternary_bool

\cs_new_protected:Npn \__numodel_parse_ternary:nN #1 #2
  {
    \bool_set_false:N \l__numodel_is_ternary_bool
    \regex_match:nnT { (.+) \? (.+) \: (.+) } {#1}
      {
        \bool_set_true:N \l__numodel_is_ternary_bool
        \tl_set:Nn \l__numodel_cond_tl {#1}
        \regex_replace_once:nnN { \s*(.+?) \s* \? .+ } { \1 } \l__numodel_cond_tl
        \tl_set:Nn \l__numodel_true_tl {#1}
        \regex_replace_once:nnN { .+? \? \s* (.+?) \s* \: .+ } { \1 } \l__numodel_true_tl
        \tl_set:Nn \l__numodel_false_tl {#1}
        \regex_replace_once:nnN { .+ \: \s* (.+?) \s* \Z } { \1 } \l__numodel_false_tl
        \__numodel_vars_to_display:N \l__numodel_cond_tl
        \__numodel_vars_to_display:N \l__numodel_true_tl
        \__numodel_vars_to_display:N \l__numodel_false_tl
      }
  }

% ====================================================================
% \mrule[keys]{varname}{calculation}
% ====================================================================
% Declare a model rule.
%
%   #1 (star)      -- starred variant: multiline IF/THEN/ELSE/ENDIF
%                     for ternary expressions
%   #2 (optional)  -- keys: alias, aliasleft, aliasright
%                       alias       -- replaces the whole display row
%                       aliasleft   -- replaces the left-hand side
%                                      (symbol)
%                       aliasright  -- replaces the right-hand side
%                                      (calculation)
%                     Use \cdots for omitted parts.
%   #3             -- name (string) of the left-hand variable
%   #4             -- calculation (with \modXxx macros and \fpeval
%                     syntax)
%
% With alias or aliasright the ternary logic is skipped (display
% shows the alias content verbatim, not IF/THEN/ELSE).  Execution
% always uses the calculation from #4.
%
% The calculation accepts all \fpeval operators:
%   +, -, *, /, ^, sign(), abs(), sin(), cos(), sqrt(), round()
%   Ternary: cond ? expr_true : expr_false
%   Logical: && (AND), || (OR), comparisons (<, >, <=, >=)
%
% The rule is stored for both display (\textmodel) and execution
% (\computemodel).
%
% Examples:
%   \mrule{modFres}{\modM * \modA}
%   \mrule{modFw}{sign(\modV) * \modK * \modV^2}
%   \mrule{modA}{(\modT < \modTq) || (\modT > \modTdq) ? \modA + \modDa : \modA - \modDa}
%   \mrule[aliasright=\cdots]{modT}{\modT + \modDt}
%   \mrule[aliasleft=a_x]{modAx}{\modFgx / \modM}
%   \mrule[alias={\text{(hidden)}}]{modAy}{\modFgy / \modM}

\NewDocumentCommand{\mrule}{ s O{} m m }{%
  \typeout{MRULE~start:~#3}
  \bool_set_false:N \l__numodel_is_ternary_bool
  % Parse keys (prefix, alias, aliasleft, aliasright; gridx/gridy ignored)
  \tl_clear:N \l__numodel_alias_tl
  \tl_clear:N \l__numodel_aliasleft_tl
  \tl_clear:N \l__numodel_aliasright_tl
  \tl_clear:N \l__numodel_cmd_prefix_tl
  \keys_set:nn { numodel / mvar } {#2}
  \__numodel_resolve_eff_prefix:
  \__numodel_with_prefix:Nn \l__numodel_eff_prefix_tl
    { \__numodel_mrule_body:nnn {#1}{#3}{#4} }%
  \typeout{MRULE~done:~#3}
}

\cs_new_protected:Npn \__numodel_mrule_body:nnn #1 #2 #3
  {
    \int_gincr:N \g_mrule_counter_int
    \tl_set:Ne \l__numodel_fullname_tl { \g_numodel_current_prefix_tl #2 }
    % Store execution data (target = full name)
    \seq_gput_right:Nx \g_mrule_calc_seq
      { { \l__numodel_fullname_tl } { \exp_not:n {#3} } }
    % Left-hand side: aliasleft or default display name
    \tl_if_blank:VTF \l__numodel_aliasleft_tl
      { \tl_set:Ne \l__numodel_lhs_tl { \use:c { \l__numodel_fullname_tl text } } }
      { \tl_set_eq:NN \l__numodel_lhs_tl \l__numodel_aliasleft_tl }
    % Generate display
    \tl_if_blank:VTF \l__numodel_alias_tl
      {
        \tl_if_blank:VTF \l__numodel_aliasright_tl
          {
            % Automatic display generation (ternary or normal)
            \__numodel_parse_ternary:nN {#3} \l__numodel_display_tl
            \bool_if:NTF \l__numodel_is_ternary_bool
              {
                \IfBooleanTF {#1}
                  {
                    % Starred ternary -> multiline IF/THEN/ELSE/ENDIF
                    \tl_set:Ne \l__numodel_display_tl
                      {
                        \__numodel_kwt:n {if}
                        \exp_not:V \l__numodel_cond_tl
                        \__numodel_kwt:n {then_nl}
                      }
                    \seq_gput_right:NV \g_mrule_seq \l__numodel_display_tl
                    \seq_gput_right:Nn \g_mrule_type_seq { rule }
                    \tl_set:Ne \l__numodel_display_tl
                      {
                        \exp_not:n { \quad }
                        \exp_not:V \l__numodel_lhs_tl \exp_not:n { \,=\, }
                        \exp_not:V \l__numodel_true_tl
                      }
                    \seq_gput_right:NV \g_mrule_seq \l__numodel_display_tl
                    \seq_gput_right:Nn \g_mrule_type_seq { cont }
                    \seq_gput_right:Ne \g_mrule_seq { \__numodel_kwt:n {else_nl} }
                    \seq_gput_right:Nn \g_mrule_type_seq { cont }
                    \tl_set:Ne \l__numodel_display_tl
                      {
                        \exp_not:n { \quad }
                        \exp_not:V \l__numodel_lhs_tl \exp_not:n { \,=\, }
                        \exp_not:V \l__numodel_false_tl
                      }
                    \seq_gput_right:NV \g_mrule_seq \l__numodel_display_tl
                    \seq_gput_right:Nn \g_mrule_type_seq { cont }
                    \seq_gput_right:Ne \g_mrule_seq { \__numodel_kwt:n {endif_nl} }
                    \seq_gput_right:Nn \g_mrule_type_seq { cont }
                  }
                  {
                    % Unstarred ternary -> single-line
                    \tl_set:Ne \l__numodel_display_tl
                      {
                        \__numodel_kwt:n {if}
                        \exp_not:V \l__numodel_cond_tl
                        \__numodel_kwt:n {then}
                        \exp_not:V \l__numodel_lhs_tl \exp_not:n { \,=\, }
                        \exp_not:V \l__numodel_true_tl
                        \__numodel_kwt:n {else}
                        \exp_not:V \l__numodel_lhs_tl \exp_not:n { \,=\, }
                        \exp_not:V \l__numodel_false_tl
                        \__numodel_kwt:n {endif}
                      }
                  }
              }
              {
                % Ordinary assignment: lhs = rhs
                \tl_set:Nn \l__numodel_rhs_tl {#3}
                \__numodel_vars_to_display:N \l__numodel_rhs_tl
                \typeout{MRULE~rhs~after:~\tl_to_str:N \l__numodel_rhs_tl}
                \tl_set:Ne \l__numodel_display_tl
                  {
                    \exp_not:V \l__numodel_lhs_tl \exp_not:n { \,=\, }
                    \exp_not:V \l__numodel_rhs_tl
                  }
                \typeout{MRULE~display:~\tl_to_str:N \l__numodel_display_tl}
              }
          }
          {
            % aliasright: lhs = aliasright (no ternary processing)
            \tl_set:Ne \l__numodel_display_tl
              {
                \exp_not:V \l__numodel_lhs_tl \exp_not:n { \,=\, }
                \exp_not:V \l__numodel_aliasright_tl
              }
          }
      }
      {
        % alias: replace the whole row
        \tl_set:Ne \l__numodel_display_tl { \exp_not:V \l__numodel_alias_tl }
      }
    % For starred ternary everything is already pushed above
    \bool_if:nF { \l__numodel_is_ternary_bool && #1 }
      {
        \seq_gput_right:NV \g_mrule_seq \l__numodel_display_tl
        \seq_gput_right:Nn \g_mrule_type_seq { rule }
      }
    % Pillar A — Lua-side rule registration.  Detokenize the raw
    % expression so Lua sees the macro names, not the evaluated
    % fp values.  Lua handles both the dependency extraction
    % (build_deps) and the flow detection (classify_flows) at
    % \graphicmodel time.
    \tl_set:Ne \l__numodel_scratch_tl { \detokenize {#3} }
    \directlua{ numodel.add_rule(
      "\g_numodel_current_prefix_tl",
      "\l__numodel_fullname_tl",
      "\luaescapestring{\l__numodel_scratch_tl}",
      "\bool_if:NTF \l__numodel_is_ternary_bool { ternary } { calc }") }
  }

% ====================================================================
% \mruletext{free text}
% ====================================================================
% Adds a free-text row in the model table.  Useful for structure
% that does not fit as \mrule, e.g.:
%   \mruletext{\text{IF } t < T/4 \text{ THEN}}
%   \mruletext{\quad a = a + da}
%   \mruletext{\text{ENDIF}}
%
% NOT executed by \computemodel (display only).

\NewDocumentCommand{\mruletext}{ O{} m }{%
  \tl_clear:N \l__numodel_cmd_prefix_tl
  \keys_set:nn { numodel / cmd } {#1}
  \__numodel_resolve_eff_prefix:
  \__numodel_with_prefix:Nn \l__numodel_eff_prefix_tl
    {
      \int_gincr:N \g_mrule_counter_int
      \seq_gput_right:Nn \g_mrule_seq {#2}
      \seq_gput_right:Nn \g_mrule_type_seq { rule }
    }
}

% ====================================================================
% \mstop[keys]{condition}
% \mstop*[keys]{condition}   (star reserved; currently identical)
% ====================================================================
% Declare the model's stop condition.
% Display: "IF condition THEN STOP ENDIF"
% Execution: the model stops when the condition is true (evaluates to 1).
%
% The condition uses \fpeval syntax with \mvar variables:
%   \mstop{\modT >= \modTmax}
%   \mstop{\modV <= 0}
%
% Supported keys:
%   alias={...}  -- replaces the whole display row (execution
%                   remains intact)
% aliasleft/aliasright are not (yet) supported for \mstop.

\tl_new:N \l__numodel_stop_tl
\NewDocumentCommand{\mstop}{ s O{} m }{%
  \typeout{MSTOP~start}
  % Parse keys (only alias is honoured; prefix via key resolver)
  \tl_clear:N \l__numodel_alias_tl
  \tl_clear:N \l__numodel_aliasleft_tl
  \tl_clear:N \l__numodel_aliasright_tl
  \tl_clear:N \l__numodel_cmd_prefix_tl
  \keys_set:nn { numodel / mvar } {#2}
  \__numodel_resolve_eff_prefix:
  \__numodel_with_prefix:Nn \l__numodel_eff_prefix_tl
    { \__numodel_mstop_body:n {#3} }%
  \typeout{MSTOP~done}
}

\cs_new_protected:Npn \__numodel_mstop_body:n #1
  {
    \int_gincr:N \g_mrule_counter_int
    % Store expression for execution (always the actual condition)
    \tl_gset:Nn \g_numodel_stop_expr_tl {#1}
    % Generate display
    \tl_if_blank:VTF \l__numodel_alias_tl
      {
        \tl_set:Nn \l__numodel_stop_tl {#1}
        \__numodel_vars_to_display:N \l__numodel_stop_tl
        \tl_set:Ne \l__numodel_display_tl
          {
            \__numodel_kwt:n {if}
            \exp_not:V \l__numodel_stop_tl
            \__numodel_kwt:n {then_nl}
            \__numodel_kwt:n {stop}
            \__numodel_kwt:n {endif_nl}
          }
      }
      {
        \tl_set:Ne \l__numodel_display_tl { \exp_not:V \l__numodel_alias_tl }
      }
    \seq_gput_right:NV \g_mrule_seq \l__numodel_display_tl
    \seq_gput_right:Nn \g_mrule_type_seq { rule }
  }

% ====================================================================
% \textmodel
% ====================================================================
% Builds a tabular with numbered model rules on the left and initial
% values on the right.  Can be placed anywhere, e.g. in a subfigure
% next to a graphic model.

\tl_new:N  \g__numodel_table_tl
\int_new:N \l__numodel_row_int
\int_new:N \l__numodel_dispnum_int
\int_new:N \l__numodel_startrow_int
\int_new:N \l__numodel_numrules_int
\int_new:N \l__numodel_numstarts_int

% Emit one initial-value cell with alias-key support.
%   - \<name>alias not empty       -> replaces whole cell (inside $...$)
%   - \<name>aliasleft not empty   -> replaces left symbol, else
%                                     \<name>text
%   - \<name>aliasright not empty  -> replaces right value, else
%                                     \<name>qty (units=true) or
%                                     \<name>num (units=false)
\cs_new_protected:Npn \__numodel_emit_startcell:n #1
  {
    % Alias-key contents are wrapped in a TeX group { ... } before
    % being placed in the table.  Two reasons:
    %  - the bare-macro form (\exp_not:c) keeps expansion-fragile
    %    contents like \cdots or \ldots from expanding at table-build
    %    time, so amsmath internals (\extrap@, \@cdots, ...) only run
    %    later inside the proper math context;
    %  - the surrounding group bounds the \futurelet lookahead that
    %    \cdots and friends perform when they expand: without it the
    %    lookahead crosses tabularray's cell boundary into document
    %    tokens, producing "Undefined control sequence \DN@" cascades.
    \str_if_eq:eeTF { \tl_to_str:c { #1 alias } } { }
      {
        \tl_gput_right:Nn \g__numodel_table_tl { $ }
        \str_if_eq:eeTF { \tl_to_str:c { #1 aliasleft } } { }
          {
            \tl_gput_right:Nx \g__numodel_table_tl
              { { \exp_not:c { #1 text } } }
          }
          {
            \tl_gput_right:Nx \g__numodel_table_tl
              { { \exp_not:c { #1 aliasleft } } }
          }
        \tl_gput_right:Nn \g__numodel_table_tl { = }
        \str_if_eq:eeTF { \tl_to_str:c { #1 aliasright } } { }
          {
            % Emit \<name>qty or \<name>num as a token name without
            % expanding it; siunitx macros must only expand at
            % \tl_use:N time, in the correct tabular context.
            \bool_if:NTF \g__numodel_units_bool
              {
                \tl_gput_right:Nx \g__numodel_table_tl
                  { \exp_not:c { #1 qty } }
              }
              {
                \tl_gput_right:Nx \g__numodel_table_tl
                  { \exp_not:c { #1 num } }
              }
          }
          {
            \tl_gput_right:Nx \g__numodel_table_tl
              { { \exp_not:c { #1 aliasright } } }
          }
        \tl_gput_right:Nn \g__numodel_table_tl { $ }
      }
      {
        \tl_gput_right:Nn \g__numodel_table_tl { $ }
        \tl_gput_right:Nx \g__numodel_table_tl
          { { \exp_not:c { #1 alias } } }
        \tl_gput_right:Nn \g__numodel_table_tl { $ }
      }
  }

\cs_new_protected:Npn \__numodel_build_table:
  {
    \typeout{BUILD_TABLE~rules:~\seq_use:Nn \g_mrule_seq {|}}
    \typeout{BUILD_TABLE~types:~\seq_use:Nn \g_mrule_type_seq {|}}
    \int_set:Nn \l__numodel_numrules_int  { \seq_count:N \g_mrule_seq }
    \int_set:Nn \l__numodel_numstarts_int { \seq_count:N \g_mvar_start_seq }
    \typeout{BUILD_TABLE~numrules:~\int_use:N \l__numodel_numrules_int}
    \typeout{BUILD_TABLE~numstarts:~\int_use:N \l__numodel_numstarts_int}
    \int_zero:N \l__numodel_row_int
    \int_zero:N \l__numodel_dispnum_int
    \int_zero:N \l__numodel_startrow_int
    \tl_gset:Nn \g__numodel_table_tl { \begin }
    \tl_gput_right:Ne \g__numodel_table_tl
      { { \tl_use:N \g__numodel_tblrenv_tl } }
    \tl_gput_right:Nn \g__numodel_table_tl
      { [theme={numodel}]{colspec={r|l|l}, rowhead=1} & \textbf }
    \tl_gput_right:Ne \g__numodel_table_tl
      { { \__numodel_kw:n {th_model} } }
    \tl_gput_right:Nn \g__numodel_table_tl
      { & \textbf }
    \tl_gput_right:Ne \g__numodel_table_tl
      { { \__numodel_kw:n {th_initvals} } }
    \tl_gput_right:Nn \g__numodel_table_tl
      { \\ \hline }
    \typeout{BUILD_TABLE~before-step~table:~\tl_to_str:N \g__numodel_table_tl}
    \int_step_inline:nn { \l__numodel_numrules_int }
      {
        \typeout{BUILD_TABLE~iter~row:~\int_use:N \l__numodel_row_int}
        \int_incr:N \l__numodel_row_int
        \int_incr:N \l__numodel_startrow_int
        \typeout{BUILD_TABLE~row:~\int_use:N \l__numodel_row_int~type:~\seq_item:Nn \g_mrule_type_seq { \l__numodel_row_int }}
        % Check type: rule -> show number, cont -> blank
        \str_if_eq:eeTF
          { \seq_item:Nn \g_mrule_type_seq { \l__numodel_row_int } }
          { cont }
          {
            % Continuation rows: no number
            \tl_gput_right:Ne \g__numodel_table_tl
              {
                \exp_not:n { \rule{0pt}{2.6ex} }
                \exp_not:n { & $ }
                \exp_not:e { \seq_item:Nn \g_mrule_seq { \l__numodel_row_int } }
                \exp_not:n { $ & }
              }
          }
          {
            % Numbered row
            \int_incr:N \l__numodel_dispnum_int
            \typeout{BUILD_TABLE~emit~rule~num:~\int_use:N \l__numodel_dispnum_int~content:~\seq_item:Nn \g_mrule_seq { \l__numodel_row_int }}
            \tl_gput_right:Ne \g__numodel_table_tl
              {
                \exp_not:n { \rule{0pt}{2.6ex} }
                \int_use:N \l__numodel_dispnum_int
                \exp_not:n { & $ }
                \exp_not:e { \seq_item:Nn \g_mrule_seq { \l__numodel_row_int } }
                \exp_not:n { $ & }
              }
            \typeout{BUILD_TABLE~after-emit~table~tail:~\tl_tail:N \g__numodel_table_tl}
          }
        % Initial-value column (independent of rule/cont)
        \int_compare:nNnTF { \l__numodel_startrow_int } > { \l__numodel_numstarts_int }
          {
            \tl_gput_right:Nn \g__numodel_table_tl { \\ }
          }
          {
            \exp_args:Ne \__numodel_emit_startcell:n
              { \seq_item:Nn \g_mvar_start_seq { \l__numodel_startrow_int } }
            \tl_gput_right:Nn \g__numodel_table_tl { \\ }
          }
      }
    \int_while_do:nNnn { \l__numodel_startrow_int } < { \l__numodel_numstarts_int }
      {
        \int_incr:N \l__numodel_startrow_int
        \tl_gput_right:Nn \g__numodel_table_tl { & & }
        \exp_args:Ne \__numodel_emit_startcell:n
          { \seq_item:Nn \g_mvar_start_seq { \l__numodel_startrow_int } }
        \tl_gput_right:Nn \g__numodel_table_tl { \\ }
      }
    \tl_gput_right:Nn \g__numodel_table_tl { \end }
    \tl_gput_right:Ne \g__numodel_table_tl
      { { \tl_use:N \g__numodel_tblrenv_tl } }
  }

\NewDocumentCommand{\textmodel}{ O{} }{%
  \typeout{TEXTMODEL~start}
  \tl_clear:N \l__numodel_cmd_prefix_tl
  \tl_clear:N \l__numodel_cmd_units_tl
  \tl_clear:N \l__numodel_cmd_tblrenv_tl
  \keys_set:nn { numodel / cmd } {#1}
  \__numodel_resolve_eff_prefix:
  % Temporary units override: save the global value, apply the key
  % value before build, restore afterwards.  This way
  % \textmodel[units=false] flips one render without touching the
  % global \numodelsetup state.
  \bool_set_eq:NN \l__numodel_saved_units_bool \g__numodel_units_bool
  \tl_if_empty:NF \l__numodel_cmd_units_tl
    {
      \str_if_eq:VnTF \l__numodel_cmd_units_tl { true }
        { \bool_gset_true:N  \g__numodel_units_bool }
        { \bool_gset_false:N \g__numodel_units_bool }
    }
  % Same dance for tblrenv: save, apply per-call override, restore.
  \tl_set_eq:NN \l__numodel_saved_tblrenv_tl \g__numodel_tblrenv_tl
  \tl_if_empty:NF \l__numodel_cmd_tblrenv_tl
    { \tl_gset_eq:NN \g__numodel_tblrenv_tl \l__numodel_cmd_tblrenv_tl }
  \__numodel_with_prefix:Nn \l__numodel_eff_prefix_tl
    {
      \group_begin:
        \__numodel_apply_dsep:
        \__numodel_build_table:
        \typeout{TEXTMODEL~table:~\tl_to_str:N \g__numodel_table_tl}
        \tl_use:N \g__numodel_table_tl
      \group_end:
    }%
  \bool_gset_eq:NN \g__numodel_units_bool \l__numodel_saved_units_bool
  \tl_gset_eq:NN \g__numodel_tblrenv_tl \l__numodel_saved_tblrenv_tl
  \typeout{TEXTMODEL~done}
}
\bool_new:N \l__numodel_saved_units_bool
\tl_new:N \l__numodel_saved_tblrenv_tl

% ====================================================================
% \computemodel
% ====================================================================
% Run the model with the Euler method:
%   1. Record all variable values in Lua
%   2. Check stop condition (\mstop)
%   3. Execute every \mrule rule (in declaration order)
%   4. Repeat until stop or safety limit (default 20000 iterations)
%
% Requires: at least one \mstop and at least one \mrule.
% Performance: ~440 steps in ~4s, ~20000 steps in ~60s.

\cs_new_protected:Npn \__numodel_exec_rule:nn #1 #2
  {
    \cs_gset:cpe {#1} { \fp_eval:n {#2} }
  }

% Record all current variable values in Lua (O(1) per variable).
\cs_new_protected:Npn \__numodel_lua_record_all:
  {
    \seq_map_inline:Nn \g_mvar_names_seq
      {
        \directlua{ numodel.record(
          "\g_numodel_current_prefix_tl", "##1", \use:c{##1}) }
      }
    \directlua{ numodel.end_step("\g_numodel_current_prefix_tl") }
  }

% Set min/max TeX macros from Lua data after the simulation finishes.
\cs_new_protected:Npn \__numodel_set_minmax_from_lua:
  {
    \seq_map_inline:Nn \g_mvar_names_seq
      {
        \cs_gset:cpe { ##1 min }
          { \directlua{ numodel.get_min("\g_numodel_current_prefix_tl", "##1") } }
        \cs_gset:cpe { ##1 max }
          { \directlua{ numodel.get_max("\g_numodel_current_prefix_tl", "##1") } }
      }
  }


\NewDocumentCommand{\computemodel}{ O{} }{%
  \typeout{COMPUTE~start}
  \tl_clear:N \l__numodel_cmd_prefix_tl
  \keys_set:nn { numodel / cmd } {#1}
  \__numodel_resolve_eff_prefix:
  \__numodel_with_prefix:Nn \l__numodel_eff_prefix_tl
    { \__numodel_compute_body: }%
  \typeout{COMPUTE~done}
}

\cs_new_protected:Npn \__numodel_compute_body:
  {
    \tl_if_empty:NT \g_numodel_stop_expr_tl
      { \msg_error:nn { numodel } { no-stop } }
    \int_gzero:N \g_numodel_steps_int
    % Initialise Lua storage for this prefix
    \directlua{ numodel.init("\g_numodel_current_prefix_tl") }
    % Main loop
    \bool_gset_false:N \g_tmpa_bool
    \bool_do_until:Nn \g_tmpa_bool
      {
        \__numodel_lua_record_all:
        \int_gincr:N \g_numodel_steps_int
        \int_compare:nNnT
          { \fp_eval:n { \g_numodel_stop_expr_tl } } = { 1 }
          { \bool_gset_true:N \g_tmpa_bool }
        \bool_if:NF \g_tmpa_bool
          {
            \seq_map_inline:Nn \g_mrule_calc_seq
              { \__numodel_exec_rule:nn ##1 }
          }
        \int_compare:nNnT
          { \g_numodel_steps_int } > { \g_numodel_maxiter_int }
          {
            \bool_gset_true:N \g_tmpa_bool
            \msg_warning:nn { numodel } { maxiter }
          }
      }
    \__numodel_set_minmax_from_lua:
    % Public per-prefix step-count accessor: \<prefix>steps holds the
    % number of recorded samples (= iterations of the loop above).
    \cs_gset:cpe { \g_numodel_current_prefix_tl steps }
      { \int_use:N \g_numodel_steps_int }
  }

% On-demand coordinates from Lua (after \computemodel).  Arguments
% are SHORT names; the current prefix is prepended automatically.
% See \mcoordsp for the form with an explicit prefix.
%
% Fully expandable (a plain \def around \directlua, no xparse
% wrapper).  Works directly inside \addplot coordinates
% {\mcoords{T}{V}} and pgfplots' math-expression parser, without
% pre-expansion via \edef.
\cs_new:Npn \mcoords #1#2
  {
    \directlua{ numodel.get_coords(
      "\g_numodel_current_prefix_tl",
      "\g_numodel_current_prefix_tl" .. "#1",
      "\g_numodel_current_prefix_tl" .. "#2") }
  }

% Variant with explicit prefix as first argument (instead of an
% optional argument, so the macro stays fully expandable).
\cs_new:Npn \mcoordsp #1#2#3
  {
    \directlua{ numodel.get_coords(
      "#1", "#1" .. "#2", "#1" .. "#3") }
  }

% Value of variable #1 at step #2 (0-based).  The variable name is
% SHORT (current prefix is prepended automatically).  See \mstepp for
% the form with an explicit prefix.
%
% Example -- tangent line at step 0:
%   \addplot[domain=0:\modTmax]
%     {\mstep{V}{0} + \mstep{A}{1} * (x - \mstep{T}{0})};
\cs_new:Npn \mstep #1#2
  {
    \directlua{ numodel.get_step(
      "\g_numodel_current_prefix_tl",
      "\g_numodel_current_prefix_tl" .. "#1", #2) }
  }

\cs_new:Npn \mstepp #1#2#3
  {
    \directlua{ numodel.get_step("#1", "#1" .. "#2", #3) }
  }

% ====================================================================
% \modelreset -- REMOVED
% ====================================================================
% \modelreset no longer exists.  Use \newmodelprefix{<prefix>} to set
% up a new model and \switchmodelprefix{<prefix>} to switch between
% existing ones.  Namespaces are additive -- variables from earlier
% models remain available.

% ====================================================================
% \graphicmodel -- Forrester diagram from \mvar/\mrule data
% ====================================================================
% Builds a tikzpicture with:
%   - stock, valve, aux and const nodes based on type and gridx/gridy
%   - flow arrows from valve to stock (via \flowarrow)
%   - causal arrows from the dependency graph
%
% Positioning: automatic (auto-layout) or manual via
%   \mvar[gridx=N, gridy=N]{...}.  Mixed mode is supported.
% Auto-layout places stocks+valves at gridy=0, aux at gridy=1, and
% constants at gridy=2.  Variables of type system are skipped.
% Auxiliary variables that act as inflow are placed as a valve
% (not as aux).

\tl_new:N   \g__numodel_graphic_tl
% Lua-populated cache (filled by numodel.tex_writeback at the start of
% each \graphicmodel call, consumed by the emit helpers below).
\prop_new:N \l__numodel_valve_for_prop      % aux/const -> stock (inflow)
% Secondary stocks for a shared inflow valve.  Value = comma-list of
% stock names; the renderer iterates it to draw a curved branch from
% the valve, sweeping over the primary stock into each extra target.
\prop_new:N \l__numodel_valve_extras_prop
\prop_new:N \l__numodel_outvalve_for_prop   % aux/const -> stock (outflow)
% Mirror of valve_extras_prop for shared outflow valves: a curved
% branch from each extra source stock arcs over the intermediate
% stocks into the shared valve's open end.
\prop_new:N \l__numodel_outvalve_extras_prop
\prop_new:N \l__numodel_between_valve_prop  % aux/const -> source_stock (between-flow)
\prop_new:N \l__numodel_between_target_prop % aux/const -> target_stock (between-flow)
\prop_new:N \l__numodel_stock_valve_prop    % stock -> source stock (stock-as-flow)
\prop_new:N \l__numodel_stock_phantom_valve_prop % stock -> source stock (phantom flow)
% For diagram-style=forrester|edu: separate valve gridx position next
% to the stock; the variable's own gridx/gridy keeps the natural
% position at gridy=1 (aux) or gridy=2 (constant).  vpos_y is always 0
% so it lives implicitly in the emit helpers.
\prop_new:N \l__numodel_vpos_x_prop         % varname -> valve gridx
\prop_new:N \l__numodel_vpos_y_prop         % varname -> valve gridy (used when gridmaxx wrap shifted the stock row)
\tl_new:N   \l__numodel_flows_tl            % deferred flow arrows
\tl_new:N   \l__numodel_valves_tl           % deferred valve nodes (drawn on top)
\tl_new:N   \l__numodel_late_causals_tl     % causals pointing at valves

% --- Resolved render settings (set per-\graphicmodel call) ---
% These hold the effective values (after resolving diagram-style
% defaults and explicit global/local overrides) used by the emit
% helpers and by the public \flowarrow / \flowoutarrow macros.
\tl_new:N \l__numodel_flowarrow_eff_tl     % "hollow" or "filled"
\tl_new:N \l__numodel_valve_eff_tl         % "valve" | "circle" | "edu"
\tl_new:N \l__numodel_flowcloud_global_tl  % "true" or "false" (global default)

% Resolve flowarrow-style: explicit global key wins; otherwise
% diagram-style picks the default (forrester -> hollow,
% tight|edu -> filled).
\cs_new_protected:Npn \__numodel_resolve_flowarrow:
  {
    \tl_if_empty:NTF \g__numodel_flowarrow_style_tl
      {
        \str_if_eq:VnTF \g__numodel_diagram_style_tl { forrester }
          { \tl_set:Nn \l__numodel_flowarrow_eff_tl { hollow } }
          { \tl_set:Nn \l__numodel_flowarrow_eff_tl { filled } }
      }
      { \tl_set_eq:NN \l__numodel_flowarrow_eff_tl
                       \g__numodel_flowarrow_style_tl }
  }

% Resolve valve-style: explicit global key wins; otherwise
% diagram-style picks the default (forrester -> valve,
% tight|edu -> edu).
\cs_new_protected:Npn \__numodel_resolve_valve:
  {
    \tl_if_empty:NTF \g__numodel_valve_style_tl
      {
        \str_if_eq:VnTF \g__numodel_diagram_style_tl { forrester }
          { \tl_set:Nn \l__numodel_valve_eff_tl { valve } }
          { \tl_set:Nn \l__numodel_valve_eff_tl { edu   } }
      }
      { \tl_set_eq:NN \l__numodel_valve_eff_tl
                       \g__numodel_valve_style_tl }
  }

% Resolve the global flowarrow-cloud-tip default (no per-stock
% override yet; that's added on top in \__numodel_eff_flowcloud:n).
\cs_new_protected:Npn \__numodel_resolve_flowcloud:
  {
    \tl_if_empty:NTF \g__numodel_flowcloud_tl
      {
        \str_if_eq:VnTF \g__numodel_diagram_style_tl { forrester }
          { \tl_set:Nn \l__numodel_flowcloud_global_tl { true  } }
          { \tl_set:Nn \l__numodel_flowcloud_global_tl { false } }
      }
      { \tl_set_eq:NN \l__numodel_flowcloud_global_tl
                       \g__numodel_flowcloud_tl }
  }

% Per-stock flowcloud lookup.  #1 = stock varname.  Sets
% \l__numodel_flowcloud_eff_tl to "true" or "false".  Order: per-mvar
% override on the stock (via [flowarrow-cloud-tip=...]) wins over the
% global default.
\tl_new:N \l__numodel_flowcloud_eff_tl
\cs_new_protected:Npn \__numodel_eff_flowcloud:n #1
  {
    \tl_set:Ne \l__numodel_flowcloud_eff_tl { \use:c { #1 flowcloud } }
    \tl_if_empty:NT \l__numodel_flowcloud_eff_tl
      { \tl_set_eq:NN \l__numodel_flowcloud_eff_tl
                       \l__numodel_flowcloud_global_tl }
  }

% Sets the public \nmflowbody and \nmflowtip macros that the
% \flowarrow / \flowoutarrow / \flowbetweenarrow macros expand to.
% Reads the resolved flowarrow-style.
\cs_new_protected:Npn \__numodel_apply_flowarrow_style:
  {
    \str_if_eq:VnTF \l__numodel_flowarrow_eff_tl { hollow }
      {
        \cs_gset:Npn \nmflowbody { flowpipe-hollow }
        \cs_gset:Npn \nmflowtip  { flowpipe-hollow-tip }
        \cs_gset:Npn \nmflowtipcloudin  { flowpipe-hollow-cloudin }
        \cs_gset:Npn \nmflowtipcloudout { flowpipe-hollow-cloudout }
      }
      {
        \cs_gset:Npn \nmflowbody { flowpipe-filled }
        \cs_gset:Npn \nmflowtip  { flowpipe-filled-tip }
        \cs_gset:Npn \nmflowtipcloudin  { flowpipe-filled-cloudin }
        \cs_gset:Npn \nmflowtipcloudout { flowpipe-filled-cloudout }
      }
  }

% Sets the tikz node-style name for the valve and the boolean that
% controls whether the valve carries the variable's display label.
%   valve-style=valve  -> [valve-forrester], no label
%   valve-style=circle -> [valve-circle],   no label
%   valve-style=edu    -> [valve-edu],      with label
\tl_new:N   \l__numodel_valve_node_tl
\bool_new:N \l__numodel_valve_label_bool
\cs_new_protected:Npn \__numodel_apply_valve_style:
  {
    \str_case:VnF \l__numodel_valve_eff_tl
      {
        { valve  }
          {
            \tl_set:Nn \l__numodel_valve_node_tl { valve-forrester }
            \bool_set_false:N \l__numodel_valve_label_bool
          }
        { circle }
          {
            \tl_set:Nn \l__numodel_valve_node_tl { valve-circle }
            \bool_set_false:N \l__numodel_valve_label_bool
          }
        { edu }
          {
            \tl_set:Nn \l__numodel_valve_node_tl { valve-edu }
            \bool_set_true:N \l__numodel_valve_label_bool
          }
      }
      {
        \tl_set:Nn \l__numodel_valve_node_tl { valve-edu }
        \bool_set_true:N \l__numodel_valve_label_bool
      }
  }

% --- Diagram-style mode flag ---
% True at diagram-style=forrester|edu; consulted by \__numodel_place_node
% to decide between single-emit and natural+phantom-valve double-emit.
% Set in \__numodel_build_graphic from \g__numodel_diagram_style_tl.
\bool_new:N \l__numodel_keep_natural_bool


\cs_new_protected:Npn \__numodel_build_graphic:
  {
    \typeout{BUILD:~step~0~reset~auto-positions}
    % --- Reset gridx/gridy to the original user input ---
    % Auto-layout writes positions to var.gridx/gridy.  On a second
    % \graphicmodel call those would, without reset, be treated as
    % "manually placed".  The initial values are saved in
    % gridxinit/gridyinit by \mvar -- copy them back.
    \seq_map_inline:Nn \g_mvar_names_seq
      {
        \cs_gset:cpe { ##1 gridx } { \use:c { ##1 gridxinit } }
        \cs_gset:cpe { ##1 gridy } { \use:c { ##1 gridyinit } }
      }
    % Diagram-style mode: forrester|edu keep aux/const at their natural
    % gridy with a phantom valve next to the stock; tight collapses the
    % aux/const onto the valve position.  Consumed by \__numodel_place_node
    % to dispatch between single-emit and natural+phantom emit.
    \bool_set_false:N \l__numodel_keep_natural_bool
    \str_if_eq:VnT \g__numodel_diagram_style_tl { forrester }
      { \bool_set_true:N \l__numodel_keep_natural_bool }
    \str_if_eq:VnT \g__numodel_diagram_style_tl { edu }
      { \bool_set_true:N \l__numodel_keep_natural_bool }
    \typeout{BUILD:~step~1~lua~layout~writeback}
    % Pillar A — Lua is the single source of truth for flow detection
    % and auto-layout.  The writeback clears and fills \<var>gridx/gridy
    % and the flow-props that downstream emitters (place_node,
    % flow-builders, emit_natural_and_phantom, emit_stock_valve) read.
    \__numodel_lua_layout_writeback:
    \typeout{BUILD:~step~2~tikzpicture}
    % --- Build tikzpicture ---
    % Render order matters: each segment overlays the previous one,
    % so we draw flows first (one continuous arrow per flow), then
    % the valve nodes on top (with white fill so they cover the
    % section of the arrow that lies underneath), and finally the
    % causal arrows that point at those valves.
    \tl_gclear:N \g__numodel_graphic_tl
    \tl_clear:N  \l__numodel_flows_tl
    \tl_clear:N  \l__numodel_valves_tl
    \tl_clear:N  \l__numodel_late_causals_tl
    \tl_gput_right:Nn \g__numodel_graphic_tl
      { \begin{tikzpicture}[gridscale] }
    % Nodes (stocks, aux, const, clouds; flows -> flows_tl,
    % valves -> valves_tl)
    \seq_map_inline:Nn \g_mvar_names_seq
      { \typeout{NODE:~##1} \__numodel_place_node:n {##1} }
    \typeout{BUILD:~step~5~stock-valve~nodes}
    % Stock-as-flow valve nodes
    \prop_map_inline:Nn \l__numodel_stock_valve_prop
      {
        % ##1 = stock (e.g. modH), ##2 = source stock (e.g. modV)
        % Skip entries that are not real stock-valves (e.g. __svgx keys)
        \tl_if_in:nnF {##1} { __sv }
          { \__numodel_emit_stock_valve:nn {##1} {##2} }
      }
    % Stock-as-rate phantom-valve nodes (source stock has no matching
    % outflow term, so the inflow gets a cloud at the open end)
    \prop_map_inline:Nn \l__numodel_stock_phantom_valve_prop
      {
        \tl_if_in:nnF {##1} { __sv }
          { \__numodel_emit_stock_phantom_valve:nn {##1} {##2} }
      }
    \typeout{BUILD:~step~6~flow~arrows}
    % Flow arrows (one segment each; drawn underneath the valves)
    \tl_gput_right:NV \g__numodel_graphic_tl \l__numodel_flows_tl
    \typeout{BUILD:~step~6a~valve~nodes}
    % Valve nodes (white-filled, drawn ON TOP of the flow arrows)
    \tl_gput_right:NV \g__numodel_graphic_tl \l__numodel_valves_tl
    \typeout{BUILD:~step~6b+7~causal~arrows~(Lua)}
    % A2: causal arrows via Lua — emit_causals demultiplexes on
    % tgt_is_valve and pushes to \g__numodel_graphic_tl or
    % \l__numodel_late_causals_tl respectively.  Afterwards append
    % late-causals to graphic_tl once: that covers both the pushes
    % from the flow-builders (step 6/6a) and those from emit_causals.
    \__numodel_lua_emit_causals:
    \tl_gput_right:NV \g__numodel_graphic_tl \l__numodel_late_causals_tl
    \typeout{BUILD:~step~8~done}
    \tl_gput_right:Nn \g__numodel_graphic_tl
      { \end{tikzpicture} }
  }

% --- Node placement ---
\cs_new_protected:Npn \__numodel_place_node:n #1
  {
    \bool_set_true:N \l_tmpa_bool
    \tl_set:Ne \l__numodel_scratch_tl { \use:c { #1 type } }
    \exp_args:NV \str_if_eq:nnT \l__numodel_scratch_tl { system }
      { \bool_set_false:N \l_tmpa_bool }
    \bool_if:NT \l_tmpa_bool
      {
        \int_compare:nNnT { \use:c { #1 gridx } } = { -1 }
          { \bool_set_false:N \l_tmpa_bool }
      }
    \bool_if:NT \l_tmpa_bool
      {
        \tl_set:Ne \l__numodel_tmp_tl { \use:c { #1 text } }
        % With forrester|edu, valve vars get a dual emission:
        % the natural aux/const node + a phantom valve next to the stock.
        \bool_set_false:N \l_tmpb_bool   % is this a valve var?
        \prop_if_in:NnT \l__numodel_valve_for_prop {#1}
          { \bool_set_true:N \l_tmpb_bool }
        \prop_if_in:NnT \l__numodel_outvalve_for_prop {#1}
          { \bool_set_true:N \l_tmpb_bool }
        \prop_if_in:NnT \l__numodel_between_valve_prop {#1}
          { \bool_set_true:N \l_tmpb_bool }
        \bool_lazy_and:nnTF
          { \l__numodel_keep_natural_bool } { \l_tmpb_bool }
          { \__numodel_emit_natural_and_phantom:n {#1} }
          {
            \prop_if_in:NnTF \l__numodel_valve_for_prop {#1}
              { \__numodel_emit_valve:n {#1} }
              {
                \prop_if_in:NnTF \l__numodel_outvalve_for_prop {#1}
                  { \__numodel_emit_outvalve:n {#1} }
                  {
                    \prop_if_in:NnTF \l__numodel_between_valve_prop {#1}
                      { \__numodel_emit_between_valve:n {#1} }
                      {
                        \tl_set:Ne \l__numodel_scratch_tl { \use:c { #1 type } }
                        \exp_args:NV \str_if_eq:nnT \l__numodel_scratch_tl { stock }
                          { \__numodel_emit_stock:n {#1} }
                        \exp_args:NV \str_if_eq:nnT \l__numodel_scratch_tl { aux }
                          { \__numodel_emit_aux:n {#1} }
                        \exp_args:NV \str_if_eq:nnT \l__numodel_scratch_tl { constant }
                          { \__numodel_emit_const:n {#1} }
                      }
                  }
              }
          }
      }
  }

% --- Helper: emit a coordinate node at the valve position ---
% The flow arrow starts/ends at this coordinate (or at an offset
% relative to it, for the open-end without cloud).  The visible
% white-filled valve is added later (valves_tl) at the same
% coordinates so it overlays the arrow.
% #1 = coord id   #2 = x   #3 = y
\cs_new_protected:Npn \__numodel_emit_valve_coord:nnn #1 #2 #3
  {
    \tl_gput_right:Ne \g__numodel_graphic_tl
      {
        \exp_not:N \coordinate ~
        (#1) ~ \exp_not:n{at} ~ (#2 , ~ #3) ;
      }
  }

% --- Helper: append a valve node to the deferred valves_tl ---
% Renders \node[<valve-style>] (id) at (x,y) {<label>}; later in the
% picture so the white fill overlays the flow arrow underneath.
% #1 = node id   #2 = x coord   #3 = y coord   #4 = label tl name
\cs_new_protected:Npn \__numodel_defer_valve:nnnn #1 #2 #3 #4
  {
    \tl_put_right:Ne \l__numodel_valves_tl
      {
        \exp_not:N \node ~ [ \tl_use:N \l__numodel_valve_node_tl ] ~
        (#1) ~ \exp_not:n{at} ~ (#2 , ~ #3) ~
        {
          \bool_if:NTF \l__numodel_valve_label_bool
            { \exp_not:N $ \exp_not:V #4 \exp_not:N $ }
            { }
        } ;
      }
  }

% --- Forrester/edu-mode emission: natural node + phantom valve ---
% The variable sits as an ordinary aux/const node at its own gridy.
% The phantom valve (id #1__v) sits next to the stock at gridy=0; its
% appearance follows the resolved valve-style (label included only
% when valve-style=edu).  A causal arrow from #1 -> #1__v is drawn
% (deferred to late_causals_tl so it lands on top of the white-
% filled valve).
\tl_new:N   \l__numodel_phantom_label_tl
\cs_new_protected:Npn \__numodel_emit_natural_and_phantom:n #1
  {
    % --- 1. Natural node (aux or constant) ---
    \tl_set:Ne \l__numodel_scratch_tl { \use:c { #1 type } }
    \exp_args:NV \str_if_eq:nnT \l__numodel_scratch_tl { aux }
      { \__numodel_emit_aux:n {#1} }
    \exp_args:NV \str_if_eq:nnT \l__numodel_scratch_tl { constant }
      { \__numodel_emit_const:n {#1} }
    % --- 2. Locate the phantom valve x/y position ---
    \prop_get:NnNT \l__numodel_vpos_x_prop {#1} \l__numodel_phantom_x_tl
      {
        \prop_get:NnNF \l__numodel_vpos_y_prop {#1} \l__numodel_scratch_y_tl
          { \tl_set:Nn \l__numodel_scratch_y_tl { 0 } }
        \tl_set:Ne \l__numodel_phantom_label_tl { \use:c { #1 text } }
        % --- 3. Coordinate at the phantom valve position ---
        \__numodel_emit_valve_coord:nnn { #1 __vc }
          { \tl_use:N \l__numodel_phantom_x_tl }
          { \tl_use:N \l__numodel_scratch_y_tl }
        % --- 4. Flow arrow + (optional) cloud, one segment ---
        \prop_if_in:NnTF \l__numodel_valve_for_prop {#1}
          {
            \tl_set:Ne \l__numodel_tmp_tl
              { \prop_item:Nn \l__numodel_valve_for_prop {#1} }
            \__numodel_eff_flowcloud:V \l__numodel_tmp_tl
            \str_if_eq:VnTF \l__numodel_flowcloud_eff_tl { true }
              {
                \tl_put_right:Ne \l__numodel_flows_tl
                  {
                    \exp_not:N \flowarrow [#1 __cl] {#1 __vc}
                      { \tl_use:N \l__numodel_tmp_tl }
                  }
              }
              {
                \tl_put_right:Ne \l__numodel_flows_tl
                  {
                    \exp_not:N \flowarrow {#1 __vc}
                      { \tl_use:N \l__numodel_tmp_tl }
                  }
              }
            \__numodel_emit_valve_extras:n {#1}
          }
          {
            \prop_if_in:NnTF \l__numodel_outvalve_for_prop {#1}
              {
                \tl_set:Ne \l__numodel_tmp_tl
                  { \prop_item:Nn \l__numodel_outvalve_for_prop {#1} }
                \__numodel_eff_flowcloud:V \l__numodel_tmp_tl
                \str_if_eq:VnTF \l__numodel_flowcloud_eff_tl { true }
                  {
                    \tl_put_right:Ne \l__numodel_flows_tl
                      {
                        \exp_not:N \flowoutarrow [#1 __cl]
                          { \tl_use:N \l__numodel_tmp_tl } {#1 __vc}
                      }
                  }
                  {
                    \tl_put_right:Ne \l__numodel_flows_tl
                      {
                        \exp_not:N \flowoutarrow
                          { \tl_use:N \l__numodel_tmp_tl } {#1 __vc}
                      }
                  }
                \__numodel_emit_outvalve_extras:n {#1}
              }
              {
                % between (no cloud: both ends are stocks)
                \tl_set:Ne \l__numodel_tmp_tl
                  { \prop_item:Nn \l__numodel_between_valve_prop {#1} }
                \tl_set:Ne \l__numodel_scratch_tl
                  { \prop_item:Nn \l__numodel_between_target_prop {#1} }
                \tl_put_right:Ne \l__numodel_flows_tl
                  {
                    \exp_not:N \flowbetweenarrow
                      { \tl_use:N \l__numodel_tmp_tl }
                      {#1 __vc}
                      { \tl_use:N \l__numodel_scratch_tl }
                  }
              }
          }
        % --- 5. Defer the phantom valve drawing (overlays flow) ---
        \__numodel_defer_valve:nnnn { #1 __v }
          { \tl_use:N \l__numodel_phantom_x_tl }
          { \tl_use:N \l__numodel_scratch_y_tl }
          \l__numodel_phantom_label_tl
        % --- 6. Causal arrow natural -> phantom valve (deferred) ---
        \tl_put_right:Ne \l__numodel_late_causals_tl
          {
            \exp_not:N \draw ~ [\exp_not:n{causal}] ~
            (#1) ~ \exp_not:n{to} ~ (#1 __v) ;
          }
      }
  }
\tl_new:N \l__numodel_phantom_x_tl
\cs_generate_variant:Nn \__numodel_eff_flowcloud:n { V }

\cs_new_protected:Npn \__numodel_emit_stock:n #1
  {
    \tl_gput_right:Ne \g__numodel_graphic_tl
      {
        \exp_not:N \node ~ [\exp_not:n{stock}] ~
        (#1) ~ \exp_not:n{at} ~
        ( \use:c{#1 gridx} , ~ \use:c{#1 gridy} ) ~
        { \exp_not:N $ \exp_not:V \l__numodel_tmp_tl \exp_not:N $ } ;
      }
  }

% Tight-mode helpers: the valve sits at the variable's own gridx/
% gridy.  We emit a coordinate node (#1__vc) early so the flow
% arrow has a positioning anchor, and defer the visible valve node
% (#1) to valves_tl so it overlays the arrow with white fill.
\cs_new_protected:Npn \__numodel_emit_tight_valve_setup:n #1
  {
    \__numodel_emit_valve_coord:nnn { #1 __vc }
      { \use:c{#1 gridx} } { \use:c{#1 gridy} }
    \__numodel_defer_valve:nnnn {#1}
      { \use:c{#1 gridx} } { \use:c{#1 gridy} }
      \l__numodel_tmp_tl
  }

\cs_new_protected:Npn \__numodel_emit_valve:n #1
  {
    \__numodel_emit_tight_valve_setup:n {#1}
    \tl_set:Ne \l__numodel_scratch_tl
      { \prop_item:Nn \l__numodel_valve_for_prop {#1} }
    \__numodel_eff_flowcloud:V \l__numodel_scratch_tl
    \str_if_eq:VnTF \l__numodel_flowcloud_eff_tl { true }
      {
        \tl_put_right:Ne \l__numodel_flows_tl
          {
            \exp_not:N \flowarrow [#1 __cl] {#1 __vc}
              { \tl_use:N \l__numodel_scratch_tl }
          }
      }
      {
        \tl_put_right:Ne \l__numodel_flows_tl
          {
            \exp_not:N \flowarrow {#1 __vc}
              { \tl_use:N \l__numodel_scratch_tl }
          }
      }
    \__numodel_emit_valve_extras:n {#1}
  }

% Emit one curved inflow branch per shared-inflow extra target
% (\l__numodel_valve_extras_prop[fv]).  Each branch starts at the
% valve coordinate and sweeps over the primary stock to the extra
% target's north edge.  No-op when the valve is not shared.
\cs_new_protected:Npn \__numodel_emit_valve_extras:n #1
  {
    \prop_get:NnNT \l__numodel_valve_extras_prop {#1} \l__numodel_scratch_tl
      {
        \clist_set:NV \l__numodel_scratch_clist \l__numodel_scratch_tl
        \clist_map_inline:Nn \l__numodel_scratch_clist
          {
            \tl_put_right:Nn \l__numodel_flows_tl
              { \flowarrowbent {#1 __vc} {##1} }
          }
      }
  }

% Mirror of \__numodel_emit_valve_extras:n for shared outflow valves.
% Each extra source stock gets a curved branch arcing over the
% intermediate stocks into the shared valve's open end.
\cs_new_protected:Npn \__numodel_emit_outvalve_extras:n #1
  {
    \prop_get:NnNT \l__numodel_outvalve_extras_prop
      {#1} \l__numodel_scratch_tl
      {
        \clist_set:NV \l__numodel_scratch_clist \l__numodel_scratch_tl
        \clist_map_inline:Nn \l__numodel_scratch_clist
          {
            \tl_put_right:Nn \l__numodel_flows_tl
              { \flowoutarrowbent {##1} {#1 __vc} }
          }
      }
  }
\clist_new:N \l__numodel_scratch_clist

\cs_new_protected:Npn \__numodel_emit_outvalve:n #1
  {
    \__numodel_emit_tight_valve_setup:n {#1}
    \tl_set:Ne \l__numodel_scratch_tl
      { \prop_item:Nn \l__numodel_outvalve_for_prop {#1} }
    \__numodel_eff_flowcloud:V \l__numodel_scratch_tl
    \str_if_eq:VnTF \l__numodel_flowcloud_eff_tl { true }
      {
        \tl_put_right:Ne \l__numodel_flows_tl
          {
            \exp_not:N \flowoutarrow [#1 __cl]
              { \tl_use:N \l__numodel_scratch_tl } {#1 __vc}
          }
      }
      {
        \tl_put_right:Ne \l__numodel_flows_tl
          {
            \exp_not:N \flowoutarrow
              { \tl_use:N \l__numodel_scratch_tl } {#1 __vc}
          }
      }
    \__numodel_emit_outvalve_extras:n {#1}
  }

\cs_new_protected:Npn \__numodel_emit_between_valve:n #1
  {
    \__numodel_emit_tight_valve_setup:n {#1}
    \tl_set:Ne \l__numodel_scratch_tl
      { \prop_item:Nn \l__numodel_between_valve_prop {#1} }
    \tl_set:Ne \l__numodel_tmp_tl
      { \prop_item:Nn \l__numodel_between_target_prop {#1} }
    \tl_put_right:Ne \l__numodel_flows_tl
      {
        \exp_not:N \flowbetweenarrow
          { \tl_use:N \l__numodel_scratch_tl }
          {#1 __vc}
          { \tl_use:N \l__numodel_tmp_tl }
      }
  }

\cs_new_protected:Npn \__numodel_emit_aux:n #1
  {
    \tl_gput_right:Ne \g__numodel_graphic_tl
      {
        \exp_not:N \node ~ [\exp_not:n{aux}] ~
        (#1) ~ \exp_not:n{at} ~
        ( \use:c{#1 gridx} , ~ \use:c{#1 gridy} ) ~
        { \exp_not:N $ \exp_not:V \l__numodel_tmp_tl \exp_not:N $ } ;
      }
  }

\cs_new_protected:Npn \__numodel_emit_const:n #1
  {
    \tl_gput_right:Ne \g__numodel_graphic_tl
      {
        \exp_not:N \constnode
        {#1}
        { \use:c{#1 gridx} , ~ \use:c{#1 gridy} }
        { \exp_not:N $ \exp_not:V \l__numodel_tmp_tl \exp_not:N $ }
      }
  }

% --- Stock-as-flow valve emission ---
\cs_new_protected:Npn \__numodel_emit_stock_valve:nn #1 #2
  {
    % #1 = target stock (e.g. modH), #2 = source stock (e.g. modV)
    % Pull gridx from the saved property
    \prop_get:NnNTF \l__numodel_stock_valve_prop { #1 __svgx }
      \l__numodel_scratch_tl
      {
        % Look up the wrap-shifted y-slot; default to 0 (the row that
        % the auto layout used before any gridmaxx wrap was applied).
        \prop_get:NnNF \l__numodel_stock_valve_prop { #1 __svgy }
          \l__numodel_scratch_y_tl
          { \tl_set:Nn \l__numodel_scratch_y_tl { 0 } }
        % Coordinate at the valve position (used by the flow arrow)
        \__numodel_emit_valve_coord:nnn { #1 __svc }
          { \tl_use:N \l__numodel_scratch_tl }
          { \tl_use:N \l__numodel_scratch_y_tl }
        % Flow arrow from source-stock through valve into target-stock
        % (no cloud: both ends are stocks).  The arrow runs in one
        % continuous segment; the visible valve will be drawn ON TOP
        % later via valves_tl.
        \tl_put_right:Ne \l__numodel_flows_tl
          {
            \exp_not:N \flowbetweenarrow {#2} {#1 __svc} {#1}
          }
        % Defer the visible valve node (white fill overlays arrow)
        \tl_set:Ne \l__numodel_tmp_tl { \use:c { #2 text } }
        \__numodel_defer_valve:nnnn { #1 __sv }
          { \tl_use:N \l__numodel_scratch_tl }
          { \tl_use:N \l__numodel_scratch_y_tl }
          \l__numodel_tmp_tl
        % Causal arrow source-stock -> valve.  Bend it (curved) only
        % when the source stock and the valve sit on the same row —
        % then the straight causal would coincide with the flow pipe
        % and disappear behind it.  When they are on different rows
        % (e.g. after a gridmaxx wrap) the causal is vertical-ish and
        % stays visible without a bend.  Deferred so it lands on top
        % of the white-filled valve.
        \str_if_eq:eeTF { \tl_use:N \l__numodel_scratch_y_tl }
                        { \use:c { #2 gridy } }
          {
            \tl_put_right:Ne \l__numodel_late_causals_tl
              {
                \exp_not:N \draw ~ [\exp_not:n{causal}] ~
                (#2) ~ \exp_not:n{to} ~ [\exp_not:n{bend~left=30}] ~ (#1__sv) ;
              }
          }
          {
            \tl_put_right:Ne \l__numodel_late_causals_tl
              {
                \exp_not:N \draw ~ [\exp_not:n{causal}] ~
                (#2) ~ \exp_not:n{to} ~ (#1__sv) ;
              }
          }
      }
      { }  % no gridx found: skip
  }

% --- Stock phantom-valve emission (source stock as rate factor) ---
% Source stock acts as a rate factor for the target stock without a
% matching outflow term; render with a cloud-fed inflow valve next to
% the target stock and link the source stock via a causal arrow.
\cs_new_protected:Npn \__numodel_emit_stock_phantom_valve:nn #1 #2
  {
    % #1 = target stock (e.g. ballY), #2 = source stock (e.g. ballV)
    \prop_get:NnNTF \l__numodel_stock_phantom_valve_prop { #1 __svgx }
      \l__numodel_scratch_tl
      {
        % Look up the wrap-shifted y-slot; default to 0.
        \prop_get:NnNF \l__numodel_stock_phantom_valve_prop { #1 __svgy }
          \l__numodel_scratch_y_tl
          { \tl_set:Nn \l__numodel_scratch_y_tl { 0 } }
        % Coordinate at the valve position
        \__numodel_emit_valve_coord:nnn { #1 __svc }
          { \tl_use:N \l__numodel_scratch_tl }
          { \tl_use:N \l__numodel_scratch_y_tl }
        % Flow arrow: cloud (open end) -> valve -> target stock
        \__numodel_eff_flowcloud:n {#1}
        \str_if_eq:VnTF \l__numodel_flowcloud_eff_tl { true }
          {
            \tl_put_right:Ne \l__numodel_flows_tl
              {
                \exp_not:N \flowarrow [#1 __svcl] {#1 __svc} {#1}
              }
          }
          {
            \tl_put_right:Ne \l__numodel_flows_tl
              {
                \exp_not:N \flowarrow {#1 __svc} {#1}
              }
          }
        % Defer the visible valve node, labelled with the source stock
        \tl_set:Ne \l__numodel_tmp_tl { \use:c { #2 text } }
        \__numodel_defer_valve:nnnn { #1 __sv }
          { \tl_use:N \l__numodel_scratch_tl }
          { \tl_use:N \l__numodel_scratch_y_tl }
          \l__numodel_tmp_tl
        % Causal arrow source-stock -> valve.  Bend (curved) only when
        % source and valve sit on the same row — see the equivalent
        % comment in \__numodel_emit_stock_valve:nn for the rationale.
        \str_if_eq:eeTF { \tl_use:N \l__numodel_scratch_y_tl }
                        { \use:c { #2 gridy } }
          {
            \tl_put_right:Ne \l__numodel_late_causals_tl
              {
                \exp_not:N \draw ~ [\exp_not:n{causal}] ~
                (#2) ~ \exp_not:n{to} ~ [\exp_not:n{bend~left=30}] ~ (#1__sv) ;
              }
          }
          {
            \tl_put_right:Ne \l__numodel_late_causals_tl
              {
                \exp_not:N \draw ~ [\exp_not:n{causal}] ~
                (#2) ~ \exp_not:n{to} ~ (#1__sv) ;
              }
          }
      }
      { }  % no gridx found: skip
  }

% --------------------------------------------------------------------
% Pillar A — Lua-side layout writeback.  A single directlua call
% that runs compute_layout (build_deps, classify_flows,
% populate_occupied, auto_layout, causals) and writes the resulting
% state back to the TeX-side props (\<var>gridx/gridy and the
% flow-props) that the downstream emitters (place_node,
% flow-builders, emit_natural_and_phantom, emit_stock_valve) read.
% Lua is the single source of truth for the decision logic; TeX
% only renders.
% --------------------------------------------------------------------
\cs_new_protected:Npn \__numodel_lua_layout_writeback:
  {
    % 1. Clear the props that Lua refills, so repeated \graphicmodel
    %    calls don't produce duplicate entries.
    \prop_clear:N  \l__numodel_valve_for_prop
    \prop_clear:N  \l__numodel_valve_extras_prop
    \prop_clear:N  \l__numodel_outvalve_for_prop
    \prop_clear:N  \l__numodel_outvalve_extras_prop
    \prop_clear:N  \l__numodel_between_valve_prop
    \prop_clear:N  \l__numodel_between_target_prop
    \prop_clear:N  \l__numodel_stock_valve_prop
    \prop_clear:N  \l__numodel_stock_phantom_valve_prop
    \prop_clear:N  \l__numodel_vpos_x_prop
    \prop_clear:N  \l__numodel_vpos_y_prop
    % 2. Run the Lua pipeline.
    \directlua{ numodel.compute_layout(
      "\g_numodel_current_prefix_tl",
      "\g__numodel_diagram_style_tl",
      \int_use:N \g_numodel_gridmaxx_int) }
    % 3. Write gridx/gridy and all flow-props back.
    %    tex_writeback emits a list of expl3 statements via tex.print.
    \directlua{ numodel.tex_writeback(
      "\g_numodel_current_prefix_tl") }
  }

\cs_new_protected:Npn \__numodel_lua_emit_causals:
  {
    \directlua{ numodel.emit_causals(
      "\g_numodel_current_prefix_tl") }
  }

\NewDocumentCommand{\graphicmodel}{ O{} }{%
  \tl_clear:N \l__numodel_cmd_prefix_tl
  \tl_clear:N \l__numodel_cmd_diagstyle_tl
  \tl_clear:N \l__numodel_cmd_flowarrow_tl
  \tl_clear:N \l__numodel_cmd_valve_tl
  \tl_clear:N \l__numodel_cmd_flowcloud_tl
  \keys_set:nn { numodel / cmd } {#1}
  \__numodel_resolve_eff_prefix:
  % Temporary overrides: save the global state, apply the per-call
  % key values before build, restore afterwards.  This way
  % \graphicmodel[diagram-style=forrester] flips one render without
  % touching the global \numodelsetup state.
  \tl_set_eq:NN \l__numodel_saved_diagstyle_tl \g__numodel_diagram_style_tl
  \tl_set_eq:NN \l__numodel_saved_flowarrow_tl \g__numodel_flowarrow_style_tl
  \tl_set_eq:NN \l__numodel_saved_valve_tl     \g__numodel_valve_style_tl
  \tl_set_eq:NN \l__numodel_saved_flowcloud_tl \g__numodel_flowcloud_tl
  \tl_if_empty:NF \l__numodel_cmd_diagstyle_tl
    { \tl_gset_eq:NN \g__numodel_diagram_style_tl \l__numodel_cmd_diagstyle_tl }
  \tl_if_empty:NF \l__numodel_cmd_flowarrow_tl
    { \tl_gset_eq:NN \g__numodel_flowarrow_style_tl \l__numodel_cmd_flowarrow_tl }
  \tl_if_empty:NF \l__numodel_cmd_valve_tl
    { \tl_gset_eq:NN \g__numodel_valve_style_tl \l__numodel_cmd_valve_tl }
  \tl_if_empty:NF \l__numodel_cmd_flowcloud_tl
    { \tl_gset_eq:NN \g__numodel_flowcloud_tl \l__numodel_cmd_flowcloud_tl }
  % Resolve the effective values for this render and propagate them
  % to \nmflowbody / \nmflowtip used by the flow macros.
  \__numodel_resolve_flowarrow:
  \__numodel_resolve_valve:
  \__numodel_resolve_flowcloud:
  \__numodel_apply_flowarrow_style:
  \__numodel_apply_valve_style:
  \__numodel_with_prefix:Nn \l__numodel_eff_prefix_tl
    {
      \__numodel_build_graphic:
      \tl_use:N \g__numodel_graphic_tl
    }%
  \tl_gset_eq:NN \g__numodel_diagram_style_tl    \l__numodel_saved_diagstyle_tl
  \tl_gset_eq:NN \g__numodel_flowarrow_style_tl  \l__numodel_saved_flowarrow_tl
  \tl_gset_eq:NN \g__numodel_valve_style_tl      \l__numodel_saved_valve_tl
  \tl_gset_eq:NN \g__numodel_flowcloud_tl        \l__numodel_saved_flowcloud_tl
}
\tl_new:N \l__numodel_saved_diagstyle_tl
\tl_new:N \l__numodel_saved_flowarrow_tl
\tl_new:N \l__numodel_saved_valve_tl
\tl_new:N \l__numodel_saved_flowcloud_tl

% ====================================================================
% Plot machinery
% ====================================================================
% \diagrammodel calls the pgfplots/calcplotdims infrastructure
% through internal wrappers \__numodel_calc_plot_dims: and
% \__numodel_draw_plot:n.  Those wrappers forward to \calcplotdims
% and \drawplot from the sibling package numodel-plot.
%
% The wrappers are not \cs_new_eq:NN because numodel-plot can load
% after this package; at load time \calcplotdims and \drawplot may
% not yet exist.  With a wrapper they are looked up at expansion
% time.

\cs_new:Npn  \__numodel_calc_plot_dims: { \calcplotdims }
\cs_new:Npn  \__numodel_draw_plot:n   #1 { \drawplot {#1} }

% ====================================================================
% \diagrammodel -- convenience command for model graphs
% ====================================================================
% Usage: \diagrammodel{xvar}{yvar}[extra drawplot code]{label}
%        \diagrammodel{xvar}{yvar1,yvar2,...,yvarn}[extra]{label}
%        \diagrammodel[prefix={p1,p2,...}]{xvar}{yvar1,...}[extra]{label}
%
% Builds a complete figure with model points.  Sets axis dimensions
% from min/max, axis labels from text/unit, and emits caption and
% label automatically.
%
% When the second argument is a comma-separated list, every y-variable
% whose unit matches the first one is drawn into the same diagram with
% the y-axis scaled to the joint min/max range.  Variables with a
% mismatching unit are dropped with a warning.  Every series is drawn
% as discrete model points (no connecting lines); the seven colours
% come from Okabe & Ito's colour-blind safe palette (yellow omitted)
% and are ordered so consecutive series differ in luminance, which
% keeps them distinguishable on a greyscale printout.  The legend
% lists each kept variable's display text.
%
% When the prefix key carries a comma-separated list of prefixes, the
% same set of y-variables is plotted for each prefix; the legend then
% appends the prefix in parentheses after each display text, e.g.
% "$U_3$ (ptc)".  Series order is prefix-major: every y-variable of
% the first prefix is drawn before any of the second.  The unit
% filter is applied against the first prefix's first y-variable; the
% same short names are assumed to exist (with matching units) under
% every other prefix.  x/y axis ranges reduce over the union of all
% (prefix, yvar) pairs.

\tl_new:N \l__numodel_xname_tl
\tl_new:N \l__numodel_yname_tl
\tl_new:N \l__numodel_yfull_tl
\tl_new:N \l__numodel_yunit_first_tl
\tl_new:N \l__numodel_yqty_tl
\tl_new:N \l__numodel_plotbody_tl
\clist_new:N \l__numodel_yshort_clist
\clist_new:N \l__numodel_eff_prefixes_clist
\bool_new:N  \l__numodel_multipref_bool
\int_new:N \l__numodel_yidx_int

% Colour-blind safe palette (Okabe & Ito; yellow #F0E442 omitted for
% poor white-background contrast).  Ordered so consecutive series
% differ in luminance as well as hue, keeping them distinguishable
% on a greyscale printout.
\definecolor{numodelseriesa}{HTML}{000000} % black            (L ~ 0.00)
\definecolor{numodelseriesb}{HTML}{E69F00} % orange           (L ~ 0.69)
\definecolor{numodelseriesc}{HTML}{0072B2} % dark blue        (L ~ 0.36)
\definecolor{numodelseriesd}{HTML}{56B4E9} % sky blue         (L ~ 0.74)
\definecolor{numodelseriese}{HTML}{D55E00} % vermillion       (L ~ 0.50)
\definecolor{numodelseriesf}{HTML}{CC79A7} % reddish purple   (L ~ 0.62)
\definecolor{numodelseriesg}{HTML}{009E73} % bluish green     (L ~ 0.52)

\clist_new:N \g__numodel_series_colors_clist
\clist_gset:Nn \g__numodel_series_colors_clist
  {
    numodelseriesa , numodelseriesb , numodelseriesc , numodelseriesd ,
    numodelseriese , numodelseriesf , numodelseriesg
  }

% Return the cycling colour name for series #1 (1-based), modulo 7.
\cs_new:Npn \__numodel_series_color:n #1
  {
    \clist_item:Nn \g__numodel_series_colors_clist
      { 1 + \int_mod:nn { #1 - 1 } { 7 } }
  }

\NewDocumentCommand{\diagrammodel}{ O{} m m O{} m }
  {
    \tl_clear:N \l__numodel_cmd_prefix_tl
    \keys_set:nn { numodel / cmd } {#1}
    \__numodel_resolve_eff_prefix:
    % --- Parse prefix list -----------------------------------------
    % \l__numodel_eff_prefix_tl may now carry a comma-separated list
    % of prefixes (e.g. "ptc,pw") when the user passes prefix={..}.
    % Treat it as a clist; a single prefix is just a one-item clist.
    % After this step \l__numodel_eff_prefix_tl is overwritten with
    % the first prefix, which drives the x-axis attributes and y-axis
    % label texts (the same short names are assumed to share unit and
    % display text across prefixes).
    \clist_set:NV \l__numodel_eff_prefixes_clist
                  \l__numodel_eff_prefix_tl
    \clist_get:NN \l__numodel_eff_prefixes_clist
                  \l__numodel_eff_prefix_tl
    \bool_set_false:N \l__numodel_multipref_bool
    \int_compare:nNnT
        { \clist_count:N \l__numodel_eff_prefixes_clist } > { 1 }
      { \bool_set_true:N \l__numodel_multipref_bool }
    \tl_set:Ne \l__numodel_xname_tl { \l__numodel_eff_prefix_tl #2 }
    % --- Parse comma-separated y-variable list ----------------------
    % Keep only variables whose unitraw matches the first one's; warn
    % about the rest.  The unit check uses the first prefix.
    \clist_clear:N \l__numodel_yshort_clist
    \tl_clear:N    \l__numodel_yunit_first_tl
    \int_zero:N    \l__numodel_yidx_int
    \clist_map_inline:nn {#3}
      {
        \int_incr:N \l__numodel_yidx_int
        \tl_set:Ne \l__numodel_yfull_tl { \l__numodel_eff_prefix_tl ##1 }
        \int_compare:nNnTF { \l__numodel_yidx_int } = { 1 }
          {
            \tl_set:Ne \l__numodel_yunit_first_tl
              { \use:c { \l__numodel_yfull_tl unitraw } }
            \clist_put_right:Nn \l__numodel_yshort_clist {##1}
          }
          {
            \tl_set:Ne \l_tmpa_tl
              { \use:c { \l__numodel_yfull_tl unitraw } }
            \str_if_eq:VVTF \l_tmpa_tl \l__numodel_yunit_first_tl
              { \clist_put_right:Nn \l__numodel_yshort_clist {##1} }
              {
                \msg_warning:nnee { numodel } { unit-mismatch }
                  { \tl_use:N \l__numodel_yfull_tl }
                  { \tl_use:N \l__numodel_yunit_first_tl }
              }
          }
      }
    % Canonical y-name = first kept variable (legacy single-series
    % accessor used by the caption when only one series is plotted).
    \clist_get:NN \l__numodel_yshort_clist \l_tmpa_tl
    \tl_set:Ne \l__numodel_yname_tl
      { \l__numodel_eff_prefix_tl \l_tmpa_tl }
    % --- Axis ranges -----------------------------------------------
    \def\xmin{\use:c{\l__numodel_xname_tl min}}
    \def\xmax{\use:c{\l__numodel_xname_tl max}}
    \def\xlabelqty{\use:c{\l__numodel_xname_tl text}}
    \def\xlabelunit{\use:c{\l__numodel_xname_tl unitraw}}
    % Reduce y-range over every (prefix, yvar) pair.  Use \edef +
    % \fpeval so the per-variable min/max accessors are expanded to
    % numeric literals before the fp parser sees them.
    \def\ymin{inf}
    \def\ymax{-inf}
    \clist_map_inline:Nn \l__numodel_eff_prefixes_clist
      {
        \clist_map_inline:Nn \l__numodel_yshort_clist
          {
            \tl_set:Ne \l__numodel_yfull_tl { ##1 ####1 }
            \edef\ymin{\fpeval{min(\ymin,\use:c{\l__numodel_yfull_tl min})}}
            \edef\ymax{\fpeval{max(\ymax,\use:c{\l__numodel_yfull_tl max})}}
          }
      }
    % --- y-axis label: comma-joined display texts, common unit -----
    % Texts are taken from the first prefix (yvars under other
    % prefixes are assumed to share their display text).
    \tl_clear:N \l__numodel_yqty_tl
    \int_zero:N \l__numodel_yidx_int
    \clist_map_inline:Nn \l__numodel_yshort_clist
      {
        \int_incr:N \l__numodel_yidx_int
        \int_compare:nNnF { \l__numodel_yidx_int } = { 1 }
          { \tl_put_right:Nn \l__numodel_yqty_tl { ,~ } }
        \tl_put_right:Ne \l__numodel_yqty_tl
          { \use:c { \l__numodel_eff_prefix_tl ##1 text } }
      }
    \edef\ylabelqty{\tl_use:N \l__numodel_yqty_tl}
    \edef\ylabelunit{\tl_use:N \l__numodel_yunit_first_tl}
    \begin{figure}[H]
    \centering
    \group_begin:
    \__numodel_apply_dsep:
    % --- Build the plot body ---------------------------------------
    % Single-series legacy path triggers iff there is exactly one
    % prefix AND exactly one kept y-variable.  The two-count sum is
    % 2 iff both lists hold one item.
    \int_compare:nNnTF
        {
          \clist_count:N \l__numodel_yshort_clist +
          \clist_count:N \l__numodel_eff_prefixes_clist
        } = { 2 }
      {
        % Single-series: preserve the legacy mark-only rendering.
        \edef\dm@coords
          { \mcoordsp{\l__numodel_eff_prefix_tl}{#2}{\tl_use:N \l_tmpa_tl} }
        \__numodel_draw_plot:n
          {
            \addplot[only~marks,~mark=*,~mark~size=0.5pt]~coordinates~{\dm@coords};
            \addlegendentry{model~point}
            #4
          }
      }
      {
        % Multi-series: one mark-only \addplot per (prefix, yvar)
        % pair with cycling colour-blind safe colour and a text
        % legend.  The legend's mark is enlarged so the colour stays
        % visible there.  Loop order is prefix-major: all y-vars of
        % the first prefix come before any of the second.  When more
        % than one prefix is plotted the prefix is appended in
        % parentheses after each display text.
        \tl_clear:N \l__numodel_plotbody_tl
        \int_zero:N \l__numodel_yidx_int
        \clist_map_inline:Nn \l__numodel_eff_prefixes_clist
          {
            \clist_map_inline:Nn \l__numodel_yshort_clist
              {
                \int_incr:N \l__numodel_yidx_int
                \tl_put_right:Ne \l__numodel_plotbody_tl
                  {
                    \exp_not:N \addplot
                      [ only~marks ,~mark=* ,~mark~size=0.5pt ,
                        color = \__numodel_series_color:n
                                  { \int_use:N \l__numodel_yidx_int } ,
                        legend~image~post~style = { mark~size=2pt } ]
                    coordinates {
                      \mcoordsp { ##1 } { #2 } { ####1 }
                    } ;
                    \exp_not:N \addlegendentry
                      {
                        $ \use:c { ##1 ####1 text } $
                        \bool_if:NT \l__numodel_multipref_bool
                          { ~ ( ##1 ) }
                      }
                  }
              }
          }
        \tl_put_right:Nn \l__numodel_plotbody_tl {#4}
        \exp_args:NV \__numodel_draw_plot:n \l__numodel_plotbody_tl
      }
    \group_end:
    \caption{$\tl_use:N \l__numodel_yqty_tl$($\use:c{\l__numodel_xname_tl text}$)\text{-diagram}}
    % If the project-specific worksheet system is loaded, pick the
    % '-ws' label inside a worksheet; otherwise always the bare label.
    \cs_if_exist:NTF \ifinworksheet
      { \ifinworksheet{\label{fig:#5-ws}}{\label{fig:#5}} }
      { \label{fig:#5} }
    \end{figure}
  }

\ExplSyntaxOff

% ============================================================
% GRAPHIC MODEL - STYLES AND MACROS
% ============================================================
%
% Forrester-diagram conventions:
%
%   Stock              [stock]    Rectangle              v, s, T, x
%   Valve / inflow     [valve]    Circle on a thick      a, dv, ds
%                                 flow pipe
%   Auxiliary          [aux]      Plain circle           F_res, F_air
%   Constant           \constnode Circle + dashes        m, k, g
%
% Use \dgridx and \dgridy as the grid spacings for consistent positioning.
% ============================================================

% === Grid spacing -- defaults set via \numodelsetup {graphscalex, graphscaley} ===

% === Styles ===
\tikzset{
	% Scale coordinates to the grid spacing
	gridscale/.style={x=\dgridx cm, y=\dgridy cm},
	% Stock (state variable): rectangle
	stock/.style={
		rectangle, draw, thick,
    minimum width=3 em,
    minimum height=2 em,
		font=\small
	},
	% Auxiliary (intermediate variable): plain circle
	aux/.style={
		circle, draw, thick,
		inner sep=2pt,
    minimum size=2 em,
		font=\scriptsize
	},
	% Constant: circle (dashes drawn by the \constnode macro)
	const/.style={
		circle, draw, thick,
		inner sep=2pt,
    minimum size=2 em,
		font=\scriptsize
	},
	% Valve (legacy alias): circle on the flow pipe.  Kept for
	% backwards compatibility with user code that calls \flowarrow
	% manually.  The package's emit helpers now pick one of the
	% style variants below based on \numodelsetup{valve-style=...}.
	valve/.style={
		circle, draw, thick, fill=white,
		minimum size=2 em,
		font=\small
	},
	% Valve variant: Forrester valve symbol -- two filled triangles
	% (white inside) whose tips meet at the centre.  The symmetry
	% axis stands perpendicular to the (horizontal) flow.  The
	% bounding box is tall and narrow so the apex angle at the
	% centre is roughly 60 degrees (sharper than the 90 degrees of
	% a square box), and the horizontal "ribs" along the top and
	% bottom edges close the two triangles.
	valve-forrester/.style={
		shape=rectangle,
		draw=none, fill=none,
		inner sep=0pt, outer sep=0pt,
		%minimum width=\dgridx*0.18 cm,
    minimum width=1 em,
		%minimum height=\dgridy*0.30 cm,
    minimum height=2 em,
		path picture={%
			\draw[thick, fill=white]
				(path picture bounding box.north west) --
				(path picture bounding box.north east) --
				(path picture bounding box.center) --
				cycle;
			\draw[thick, fill=white]
				(path picture bounding box.south west) --
				(path picture bounding box.south east) --
				(path picture bounding box.center) --
				cycle;
		}
	},
	% Valve variant: empty circle (no label).
	valve-circle/.style={
		circle, draw, thick, fill=white,
		minimum size=2 em,
		font=\small
	},
	% Valve variant: circle with the variable's display label inside.
	valve-edu/.style={
		circle, draw, thick, fill=white,
		minimum size=2 em,
		font=\small
	},
	% Flow pipe variants (body without arrow tip, used for the
	% segment from open-end / cloud to the valve):
	flowpipe/.style={line width=3pt},                     % legacy alias
	flowpipe-filled/.style={line width=3pt},
	flowpipe-hollow/.style={
		line width=1pt,
		double=white, double distance=3pt
	},
	% Flow pipe variants WITH arrow tip (used for the segment from
	% valve to stock, or the open-end side of an outflow).  Both
	% variants use the Stealth arrow tip.  The hollow form
	% uses the double-line magic length 0pt-3-0 so that the arrow
	% tip merges cleanly with the gap between the two strokes.
	flowpipe-filled-tip/.style={line width=3pt, arrows={-Stealth[inset=0pt, angle=30:1em]}},
	flowpipe-hollow-tip/.style={
		line width=1pt,
		double=white, double distance=3pt,
		arrows={-Stealth[inset=0pt, angle=30:1em]}
	},
	% Flow pipe variants with the Cloud arrow tip on the open end:
	% inflow keeps the Stealth arrow head at the stock side.
	flowpipe-filled-cloudin/.style={
		line width=3pt,
		arrows={Cloud[fill=white, line width=0.8pt]-Stealth[inset=0pt, angle=30:1em]}
	},
	flowpipe-hollow-cloudin/.style={
		line width=1pt,
		double=white, double distance=3pt,
		arrows={Cloud[fill=white, line width=0.8pt]-Stealth[inset=0pt, angle=30:1em]}
	},
	% Outflow with cloud at the open end: first a Stealth tip at the
	% line end, then the cloud behind it (in arrows.meta spec, list
	% multiple tips separated by a dot so the shaft stops at the
	% Stealth tip and does not extend into the cloud; order = from
	% line outwards, hence Stealth.Cloud).
	flowpipe-filled-cloudout/.style={
		line width=3pt,
		arrows={-{Stealth[inset=0pt, angle=30:1em].Cloud[fill=white, line width=0.8pt]}}
	},
	flowpipe-hollow-cloudout/.style={
		line width=1pt,
		double=white, double distance=3pt,
		arrows={-{Stealth[inset=0pt, angle=30:1em].Cloud[fill=white, line width=0.8pt]}}
	},
	% Causal arrow (thin arrow for dependencies)
	causal/.style={thick, arrows={-Stealth[length=0.5em]}}
}

% Default values for the body / tip macros so a manual call to
% \flowarrow (outside \graphicmodel) keeps working.  The
% \graphicmodel pipeline overwrites these per render.
\providecommand{\nmflowbody}{flowpipe-filled}
\providecommand{\nmflowtip}{flowpipe-filled-tip}
\providecommand{\nmflowtipcloudin}{flowpipe-filled-cloudin}
\providecommand{\nmflowtipcloudout}{flowpipe-filled-cloudout}

% === Macro: constant node with dashes ===
% Usage: \constnode{name}{(x,y)}{$label$}
\newcommand{\constnode}[3]{%
	\node[const] (#1) at (#2) {#3};%
	\draw[thick] ([xshift=-3mm]#1.west) -- (#1.west);%
	\draw[thick] (#1.east) -- ([xshift=3mm]#1.east);%
}

% === Macro: flow pipe with valve (inflow) ===
% Draws ONE continuous flow arrow ending at the stock.  The starting
% position is the same regardless of cloud presence; passing any
% non-blank optional argument switches to the cloud-tipped variant
% so the arrow's tail becomes a white-filled cloud (drawn as a real
% PGF arrow tip, not a separate node — moves with the line).  The
% valve overlays the middle of the arrow as a separate white-filled
% node drawn later by the build pipeline.
% Style keys: \nmflowtip / \nmflowtipcloudin.
% Usage: \flowarrow[<any>]{<valve-coord>}{<stock-node>}
\NewDocumentCommand{\flowarrow}{ O{} m m }{%
	\IfBlankTF{#1}{%
		\draw[\nmflowtip] ([xshift=-3em]#2) -- ([xshift=0.5mm]#3.west);%
	}{%
		\draw[\nmflowtipcloudin] ([xshift=-3em]#2) -- ([xshift=0.5mm]#3.west);%
	}%
}

% === Macro: flow pipe with valve (outflow) ===
% Mirror of \flowarrow.  The arrow always ends at the same offset
% past the valve coordinate; passing a non-blank optional argument
% switches to the cloud-tipped variant (cloud at the open end).
% Usage: \flowoutarrow[<any>]{<stock-node>}{<valve-coord>}
\NewDocumentCommand{\flowoutarrow}{ O{} m m }{%
	\IfBlankTF{#1}{%
		\draw[\nmflowtip] (#2.east) -- ([xshift=4em]#3);%
	}{%
		\draw[\nmflowtipcloudout] (#2.east) -- ([xshift=5em]#3);%
	}%
}

% === Macro: flow pipe between two stocks ===
% One continuous arrow from source stock through valve coord into
% target stock.  No cloud option (no open end exists when both
% sides are stocks).
% Usage: \flowbetweenarrow{<source-stock>}{<valve-coord>}{<target-stock>}
\newcommand{\flowbetweenarrow}[3]{%
	\draw[\nmflowtip] (#1.east) -- ([xshift=0.5mm]#3.west);%
}

% === Macro: curved inflow branch for a shared valve ===
% Used when one inflow valve feeds more than one stock.  The primary
% target is drawn with the standard straight \flowarrow; every
% secondary target is drawn with this macro, which arcs UNDER the
% intervening stock(s).  Stocks sit on the bottom row of the
% Forrester diagram, so the area below is empty -- routing the
% branch below avoids collisions with the aux/const rows and the
% causal arrows that live above.  Both endpoints are passed
% unanchored so TikZ picks the border-intersection that best
% matches the curve's tangent at each end.
% Usage: \flowarrowbent{<valve-coord>}{<target-stock>}
\newcommand{\flowarrowbent}[2]{%
	\draw[\nmflowtip] (#1) to[bend right=30] (#2);%
}

% === Macro: curved outflow branch for a shared valve ===
% Mirror of \flowarrowbent.  Used when one outflow valve drains more
% than one stock.  Routes the branch BELOW the intervening stocks
% (same rationale as the inflow variant: nothing else lives below
% the stocks row).  Endpoints unanchored so TikZ chooses the
% cleanest border intersection.
% Usage: \flowoutarrowbent{<source-stock>}{<valve-coord>}
\newcommand{\flowoutarrowbent}[2]{%
	\draw[\nmflowtip] (#1) to[bend right=30] (#2);%
}
%</package>
%    \end{macrocode}
%
% \subsection{Translation files}
%
% Each |numodel-<LANG>.def| file populates the lookup table consulted
% by |\__numodel_kw:n|.  The package loads exactly the file matching
% the current |syntax| key, via |\InputIfFileExists| (and hence
% |kpse|), so additional languages can be supplied through
% |TEXMFHOME/tex/latex/numodel/numodel-<LANG>.def| with no package
% rebuild.  The two shipped files follow.
%
% \subsubsection{\texttt{numodel-EN.def} -- English (XMILE)}
%
% XMILE v1.0 OASIS style: ALL-CAPS, case sensitive.  Sets the default
% decimal mark to |point|.
%
%    \begin{macrocode}
%<*EN>
\ProvidesExplFile{numodel-EN.def}{2026/05/26}{0.6.0}
  {numodel~keyword~table~--~English~(XMILE)}
\cs_new:Npn \__numodel_kw_EN_if:           { IF~ }
\cs_new:Npn \__numodel_kw_EN_then:         { ~THEN~ }
\cs_new:Npn \__numodel_kw_EN_then_nl:      { ~THEN }
\cs_new:Npn \__numodel_kw_EN_else:         { ~ELSE~ }
\cs_new:Npn \__numodel_kw_EN_else_nl:      { ELSE }
\cs_new:Npn \__numodel_kw_EN_endif:        { ~ENDIF }
\cs_new:Npn \__numodel_kw_EN_endif_nl:     { ENDIF }
\cs_new:Npn \__numodel_kw_EN_and:          { ~AND~ }
\cs_new:Npn \__numodel_kw_EN_or:           { ~OR~ }
\cs_new:Npn \__numodel_kw_EN_not:          { NOT~ }
\cs_new:Npn \__numodel_kw_EN_sign:         { SIGN }
\cs_new:Npn \__numodel_kw_EN_abs:          { ABS }
\cs_new:Npn \__numodel_kw_EN_sqrt:         { SQRT }
\cs_new:Npn \__numodel_kw_EN_exp:          { EXP }
\cs_new:Npn \__numodel_kw_EN_ln:           { LN }
\cs_new:Npn \__numodel_kw_EN_sin:          { SIN }
\cs_new:Npn \__numodel_kw_EN_cos:          { COS }
\cs_new:Npn \__numodel_kw_EN_tan:          { TAN }
\cs_new:Npn \__numodel_kw_EN_asin:         { ARCSIN }
\cs_new:Npn \__numodel_kw_EN_acos:         { ARCCOS }
\cs_new:Npn \__numodel_kw_EN_atan:         { ARCTAN }
\cs_new:Npn \__numodel_kw_EN_min:          { MIN }
\cs_new:Npn \__numodel_kw_EN_max:          { MAX }
\cs_new:Npn \__numodel_kw_EN_stop:         { ~STOP~ }
\cs_new:Npn \__numodel_kw_EN_th_model:     { model }
\cs_new:Npn \__numodel_kw_EN_th_initvals:  { initial~values }
\cs_new:Npn \__numodel_kw_EN_contfoot:     { Continued~on~next~page }
\cs_new:Npn \__numodel_kw_EN_conthead:     { (Continued) }
\cs_new:Npn \__numodel_kw_EN_dsep_default: { point }
%</EN>
%    \end{macrocode}
%
% \subsubsection{\texttt{numodel-NL.def} -- Dutch (CoachTaal)}
%
% CoachTaal convention.  Sets the default decimal mark to |comma|.
%
%    \begin{macrocode}
%<*NL>
\ProvidesExplFile{numodel-NL.def}{2026/05/26}{0.6.0}
  {numodel~keyword~table~--~Dutch~(CoachTaal)}
\cs_new:Npn \__numodel_kw_NL_if:           { Als~ }
\cs_new:Npn \__numodel_kw_NL_then:         { ~Dan~ }
\cs_new:Npn \__numodel_kw_NL_then_nl:      { ~Dan }
\cs_new:Npn \__numodel_kw_NL_else:         { ~Anders~ }
\cs_new:Npn \__numodel_kw_NL_else_nl:      { Anders }
\cs_new:Npn \__numodel_kw_NL_endif:        { ~EindAls }
\cs_new:Npn \__numodel_kw_NL_endif_nl:     { EindAls }
\cs_new:Npn \__numodel_kw_NL_and:          { ~EN~ }
\cs_new:Npn \__numodel_kw_NL_or:           { ~OF~ }
\cs_new:Npn \__numodel_kw_NL_not:          { NIET~ }
\cs_new:Npn \__numodel_kw_NL_sign:         { Teken }
\cs_new:Npn \__numodel_kw_NL_abs:          { Abs }
\cs_new:Npn \__numodel_kw_NL_sqrt:         { Sqrt }
\cs_new:Npn \__numodel_kw_NL_exp:          { Exp }
\cs_new:Npn \__numodel_kw_NL_ln:           { Ln }
\cs_new:Npn \__numodel_kw_NL_sin:          { Sin }
\cs_new:Npn \__numodel_kw_NL_cos:          { Cos }
\cs_new:Npn \__numodel_kw_NL_tan:          { Tan }
\cs_new:Npn \__numodel_kw_NL_asin:         { Arcsin }
\cs_new:Npn \__numodel_kw_NL_acos:         { Arccos }
\cs_new:Npn \__numodel_kw_NL_atan:         { Arctan }
\cs_new:Npn \__numodel_kw_NL_min:          { Min }
\cs_new:Npn \__numodel_kw_NL_max:          { Max }
\cs_new:Npn \__numodel_kw_NL_stop:         { ~Stop~ }
\cs_new:Npn \__numodel_kw_NL_th_model:     { model }
\cs_new:Npn \__numodel_kw_NL_th_initvals:  { startwaarden }
\cs_new:Npn \__numodel_kw_NL_contfoot:     { Wordt~vervolgd~op~volgende~pagina }
\cs_new:Npn \__numodel_kw_NL_conthead:     { (Vervolg) }
\cs_new:Npn \__numodel_kw_NL_dsep_default: { comma }
%</NL>
%    \end{macrocode}
%
% \Finale
\endinput
